% Standalone assembly. Edit article.tex and sections/, then run build.py.
\documentclass[11pt,a4paper]{article}
\usepackage[T1]{fontenc}
\usepackage[utf8]{inputenc}
\usepackage{lmodern}
\usepackage[margin=27mm,headheight=15pt]{geometry}
\usepackage{amsmath,amssymb,amsthm,mathtools,bm}
\usepackage{booktabs,array,longtable}
\usepackage{microtype}
\usepackage[dvipsnames]{xcolor}
\usepackage{enumitem}
\usepackage{fancyhdr}
\usepackage{hyperref}
\usepackage[nameinlink,noabbrev]{cleveref}
\hypersetup{colorlinks=true,linkcolor=MidnightBlue,citecolor=MidnightBlue,urlcolor=MidnightBlue,pdftitle={Ordered Boundary Identities and Gamma-Kernel Reconstruction},pdfauthor={ProveIt Contributors}}
\urlstyle{same}
\setlength{\emergencystretch}{3em}
\numberwithin{equation}{section}
\newtheorem{theorem}{Theorem}[section]
\newtheorem{proposition}[theorem]{Proposition}
\newtheorem{lemma}[theorem]{Lemma}
\newtheorem{corollary}[theorem]{Corollary}
\theoremstyle{definition}
\newtheorem{definition}[theorem]{Definition}
\newtheorem{example}[theorem]{Example}
\theoremstyle{remark}
\newtheorem{remark}[theorem]{Remark}
\DeclareMathOperator{\Li}{Li}
\DeclareMathOperator{\FP}{FP}
\DeclareMathOperator{\Res}{Res}
\DeclareMathOperator{\CT}{CT}
\DeclareMathOperator{\supp}{supp}
\newcommand{\C}{\mathbb C}
\newcommand{\Q}{\mathbb Q}
\newcommand{\Z}{\mathbb Z}
\newcommand{\N}{\mathbb N}
\newcommand{\dd}{\,\mathrm d}
\newcommand{\Rea}{\operatorname{Re}}
\newcommand{\Ord}{\mathcal O}
\newcommand{\Kern}{\mathcal K}
\newcommand{\Def}{\mathcal D}
\newcommand{\Bell}{\mathcal B}
\newcommand{\ind}{\mathbf 1}
\newcommand{\repo}{\texttt{e3d7ec5d3f9be259c1f91edaea4389f4e79b5b03}}
\pagestyle{fancy}
\fancyhf{}
\fancyhead[L]{\small Ordered boundary identities}
\fancyhead[R]{\small ProveIt research continuation}
\fancyfoot[C]{\thepage}
\setlist{itemsep=3pt,topsep=5pt}
\begin{document}
\hypersetup{pageanchor=false}
\begin{titlepage}
\centering
\vspace*{16mm}
{\small\sffamily\color{MidnightBlue} PROVEIT \quad / \quad POLYLOGARITHMS AND THEIR ARITHMETIC BRIDGES\par}
\vspace{15mm}
{\LARGE\bfseries Ordered Boundary Identities\par}
\vspace{2mm}
{\Large\bfseries and Gamma-Kernel Reconstruction\par}
\vspace{4mm}
{\Large Boundary localization, polylogarithmic cancellation,\par
 generalized Stieltjes jets, and Gamma kernels\par}
\vspace{13mm}
{\large ProveIt Contributors\par}
\vspace{3mm}
{Prepared for Vladimir Reshetnikov\par}
\vspace{5mm}
{11 October 2026\par}
\vfill
\begin{minipage}{0.88\textwidth}
\small
\textbf{Research status.} This article supplies ordinary mathematical proofs,
exact finite algebra, and separately labelled numerical diagnostics.
It is a continuation of the current manuscript and all seven incoming
polylogarithm packages inspected at the repository revision printed below.
Historical priority is not asserted for the new formulations.
No proof-assistant formalization is claimed.
\par\medskip
\textbf{Source revision.}\quad\texttt{e3d7ec5d3f9be259c1f91edaea4389f4e79b5b03}
\end{minipage}
\vspace{10mm}
\end{titlepage}
\hypersetup{pageanchor=true}
% Begin included source: sections/01_scope
\begin{abstract}
An ordered assignment of independently shifted finite-difference weights
contains information that its determinantal symmetrization erases. We
recover that information explicitly. For two particles, two finite
summations by parts reduce the ordered double sum to a single two-shift
zeta kernel. This gives its joint meromorphic continuation, every polar
coefficient, all convergent positive-integer-order evaluations and their
finite parts, a divided-difference
digamma generating function, and a Gamma primitive. Oblique poles explain
direction-dependent values at parameters where every pointwise weight
vanishes. A cutoff generating function is then proved on an entire
two-order parameter domain: its sole depth-two polylogarithm has the
factor $(1-z)^2$, which explains the endpoint reduction. Beta
parametrization yields ordinary convergent identities involving every
generalized Stieltjes constant. At arbitrary particle number, transferring
the stencils to the ordered Vandermonde square localizes the sum to
collision blocks; an explicit finite algorithm uses at most
$\lfloor(r+1)/2\rfloor$ free summation indices. A separate Gamma-integral
basis solves the fixed-gap null equation for unequal even harmonic
tails. Its exact normalization and nonzero determinant asymptotic prove
eventual optimal reconstruction on every fixed-gap ray; the first
relative correction is an explicit harmonic-number expression.
All formulas retain their convergence domains, confluence conventions,
and finite-part coordinates. Exact algebraic tests and numerical checks
accompany the proofs.
\end{abstract}

\section{The source questions and the mathematical setting}
\label{sec:scope}

\subsection{A continuation of the current material}

The source snapshot is the ProveIt revision \repo, obtained on
11 October 2026. Both requested directories were inspected: the canonical
manuscript and the seven then-present incoming archives. The manuscript's
claim ledger remains the authority for accepted canonical results
\cite{repo:manuscript,repo:ledger}; an incoming report's delivery status
does not by itself amount to canonical proof acceptance.

The most direct target is the incoming \emph{Finite-Difference
Determinants} report \cite{repo:finite}. Its section entitled
``Unsymmetrized mixed-order sums'' asks which additional coordinates are
needed for an individual ordering of independent spectral labels.
Its existing polarization theorem gives only the sum over assignments.
It also asks for a finite description of noninteger-order cutoff
generators and for extensions to other annihilating stencils. These
questions are mathematically related: shifting the finite differences
onto the ordered summation domain exposes the boundaries at which
the extra coordinates live.

The present two-particle result answers the individual-ordering question
on an open complex domain, not only at isolated integer values.
It shows that a one-index kernel is enough. That kernel has explicit
Hurwitz expansions and, at positive integer orders, a finite polygamma
evaluation. The all-particle result is different: it gives a finite
localization formula and a bound on the number of free indices, rather
than incorrectly extrapolating a universal single-zeta formula.

A second target comes from the harmonic-tail Gamma continuation
\cite{repo-gamma}. Its exact fixed-gap reduction leaves an
$m$-dimensional null equation, where the two tail orders differ by
$2m$. Section~\ref{gb:section} constructs a basis by ordinary integrals
against $|\Gamma(\tfrac14+\tfrac{ix}{2})|^{4m}$. The distinct
decay scales of these integrals produce a nonzero limiting Vandermonde
determinant. This proves that, for every fixed $m$, the first $2b$
moments form the optimal formal prefix for all sufficiently large
$(p,q)=(2b-2m,2b)$. The theorem is complementary to the existing
eventual result with the smaller order fixed \cite{repo:quartic}. Neither assertion alone
settles every unequal pair.

The current sources also make clear why old open labels cannot simply
be reused. Several finite-order harmonic-tail reconstruction cases have
already received proofs in recent reports. The short Gaussian $S_6$
and revised $S_8$ evaluations, by contrast, remain conjectural in the
current canonical manuscript \cite{repo:ledger,repo:gaussian}. None of
the ordered identities below is represented as a proof of either
Gaussian evaluation.

\subsection{Classical ingredients and the scope of originality}

Finite summation by parts, finite differences, partial fractions,
Hurwitz continuation, and Gamma differentiation are classical.
The general rational-summability role of discrete residues is treated
by Chen--Singer and Arreche--Sitaula \cite{chen-singer,arreche-sitaula}.
Multivariable zeta functions can have intersecting polar hyperplanes
and direction-sensitive Laurent behavior; this phenomenon is not new
in itself \cite{matsumoto-ono-wakabayashi}. Beta parametrization is also
an established method for zeta products \cite{shpot-paris}.

The contributions claimed here are the explicit ordered boundary
identity, its complete consequences in this family, the finite cutoff
polylogarithm formula and its endpoint cancellation, and the
all-particle collision-block reduction. Their derivations were checked
against the inspected sources. The literature search does not establish
worldwide historical priority. In particular, finite polygamma
reduction of a rational one-index summand is a classical step; it is
the preceding reduction of the ordered ensemble that supplies the
new application.

The Gamma basis, its normalization, and its fixed-gap consequences form
the second group of contributions. Classical Fourier transforms of
Gamma products provide useful context \cite{ramanujan-transform};
the result here concerns a specific higher-power weight, its
sinh/tanh modes, and the repository's exact reconstruction determinant.

\subsection{Notation and analytic conventions}

For $\Rea a>0$ we use the principal logarithm in
$(n+a)^{-s}=\exp(-s\log(n+a))$. Thus every shifted denominator in
the original series stays in the right half-plane. Write
\[
 (s)_j=s(s+1)\cdots(s+j-1),\qquad (s)_0=1,
 \qquad \psi(a)=\frac{\Gamma'(a)}{\Gamma(a)}.
\]
For positive integers $p$,
\[
 H_n^{(p)}=\sum_{k=1}^{n}k^{-p},\qquad H_0^{(p)}=0.
\]
The same finite sum defines $H_n^{(s)}$ for complex $s$.
We use the Stieltjes convention
\begin{equation}\label{eq:stieltjes-convention}
 \zeta(1+u,a)=\frac1u+
    \sum_{m=0}^{\infty}\frac{(-1)^m\gamma_m(a)}{m!}u^m,
 \qquad \gamma_0(a)=-\psi(a).
\end{equation}
Derivatives of $\zeta$ without an argument qualifier differentiate its
first variable. Standard relations used below include
\begin{equation}\label{eq:classical-ladders}
 \partial_a\zeta(s,a)=-s\zeta(s+1,a),\qquad
 \zeta(p,a)=\frac{(-1)^p}{(p-1)!}\psi^{(p-1)}(a)\quad(p\ge2),
\end{equation}
and $\zeta'(0,a)=\log\Gamma(a)-\tfrac12\log(2\pi)$
\cite{dlmf:hurwitz,dlmf:gamma}.

An equality continued from a normally convergent domain is always an
equality of those continuations. It need not equal the result of
substituting exceptional parameters into every summand first.
Likewise, a finite part is the constant in a declared Laurent
coordinate; a hard cutoff is a different operation unless their
agreement is proved.

For $|z|<1$ our multiple-polylogarithm convention is
\begin{equation}\label{eq:mpl-convention}
 \Li_{t,s}(z,1)=\sum_{k>j\ge1}\frac{z^k}{k^t j^s}.
\end{equation}
The exponential factor makes this entire in the two order variables.
This is independent of the convergence restrictions of its value at
$z=1$.

% End included source: sections/01_scope
\tableofcontents
\clearpage
% Begin included source: sections/02_ordered
\section{The ordered two-particle boundary identity}
\label{sec:ordered}

Define
\begin{equation}\label{eq:weight-two}
 w_n(s,a)=\frac12\bigl((n+a)^{-s}-2(n+a+1)^{-s}
                           +(n+a+2)^{-s}\bigr)
\end{equation}
and retain the ordering of its labels:
\begin{equation}\label{eq:ordered-def}
 \Ord(s,t;a,b)=\sum_{0\le m<n}(n-m)^2w_m(s,a)w_n(t,b).
\end{equation}
In particular, exchanging $(s,a)$ and $(t,b)$ generally changes the
answer. Let
\begin{equation}\label{eq:kernel-def}
 \Kern(s,t;x,y)=\sum_{k=0}^{\infty}(k+x)^{-s}(k+y)^{-t},
 \qquad \Rea(s+t)>1,
\end{equation}
and set
\begin{equation}\label{eq:boundary-coefficient}
 E_s(a)=\frac{a^{-s}-(a+1)^{-s}}2.
\end{equation}

\begin{theorem}[Ordered boundary reduction]\label{thm:ordered}
For
\begin{equation}\label{eq:ordered-domain}
 \Rea a,\Rea b>0,\qquad \Rea t>1,\qquad \Rea(s+t)>0,
\end{equation}
the sum \eqref{eq:ordered-def} converges absolutely and locally
normally, and
\begin{equation}\label{eq:ordered-main}
 \boxed{\Ord(s,t;a,b)=E_s(a)\zeta(t,b+1)
 +\frac14\bigl\{\Kern(s,t;a+2,b+1)-\Kern(s,t;a,b+1)\bigr\}.}
\end{equation}
On the portion $0<\Rea(s+t)\le1$, the kernel difference means the
single combined convergent series
\begin{equation}\label{eq:kernel-difference}
 \Def(s,t;a,b)=\sum_{k=0}^{\infty}(k+b+1)^{-t}
                  \bigl\{(k+a+2)^{-s}-(k+a)^{-s}\bigr\}.
\end{equation}
All fixed finite derivatives in the four parameters can be taken
termwise on compact subsets of \eqref{eq:ordered-domain}.
\end{theorem}

\subsection{Convergence on the asymmetric domain}

Repeated integration of a second derivative gives
\begin{equation}\label{eq:second-difference-integral}
 w_n(s,a)=\frac{s(s+1)}2\int_0^1\int_0^1
               (n+a+u+v)^{-s-2}\dd u\dd v.
\end{equation}
Consequently $w_n(s,a)=O((n+1)^{-\Rea s-2})$, uniformly on compact
parameter sets, including neighborhoods of $s=0,-1$.
For a fixed outer index $n$, bound $(n-m)^2$ by $n^2$.
The inner sum of $(m+1)^{-\Rea s-2}$ is bounded when
$\Rea s>-1$, grows logarithmically at $\Rea s=-1$, and is
$O(n^{-\Rea s-1})$ when $\Rea s<-1$. Combining these bounds with
the outer weight gives convergence under \eqref{eq:ordered-domain}.
On compact sets one can decrease the real exponents slightly to
obtain a single summable majorant, so the transition at $\Rea s=-1$
causes no loss of normality. Cauchy estimates on slightly larger
compact neighborhoods justify every fixed derivative.

The first difference in \eqref{eq:kernel-difference} is
$O((k+1)^{-\Rea s-1})$. Thus its combined series has precisely the
sufficient domain $\Rea(s+t)>0$ used in the theorem. No separate
summation of divergent kernel terms is involved.

\subsection{A finite identity with both upper endpoints}

The proof works for arbitrary finite sequences before any special
function is introduced. Put $f_n=(n+a)^{-s}$ and $g_n=(n+b)^{-t}$,
and define $A_{-1}=A_0=0$ and
\begin{equation}\label{eq:A-n}
 A_n=\sum_{m=0}^{n-1}(n-m)^2(f_m-2f_{m+1}+f_{m+2}).
\end{equation}
Collecting the coefficient of each $f_k$ gives, for $n\ge1$,
\begin{equation}\label{eq:A-explicit}
 A_n=n^2f_0+(-n^2-2n+1)f_1
                  +2\sum_{k=2}^{n}f_k+f_{n+1}.
\end{equation}
Taking two backward differences yields
\begin{equation}\label{eq:A-difference}
 A_n-2A_{n-1}+A_{n-2}
          =2(f_0-f_1)-f_{n-1}+f_{n+1}\qquad(n\ge1).
\end{equation}
At $n=1$ this is a direct check, using $A_{-1}=0$.

Let $\Ord_N$ denote the restriction $n\le N$ in
\eqref{eq:ordered-def}. Shifting the three finite sums in
$4\Ord_N=\sum_{n=1}^N A_n(g_n-2g_{n+1}+g_{n+2})$ gives
\begin{align}\label{eq:finite-ordered}
 4\Ord_N={}&\sum_{n=1}^{N}g_n
              \{2(f_0-f_1)-f_{n-1}+f_{n+1}\}\notag\\
 &+g_{N+1}(A_{N-1}-2A_N)+g_{N+2}A_N.
\end{align}
This is an exact finite identity. In particular, no endpoint term
has been inferred from an asymptotic prescription.

\begin{proof}[Proof of \cref{thm:ordered}]
The upper boundary in \eqref{eq:finite-ordered} can be grouped as
\[
 A_N(g_{N+2}-g_{N+1})+g_{N+1}(A_{N-1}-A_N).
\]
Formula \eqref{eq:A-explicit} shows that $A_N$ is bounded by a sum
of $O(N^2)$ and $O(N^{1-\Rea s})$, with a logarithmic enlargement
at a borderline exponent. Its first backward difference is bounded
by the corresponding $O(N)$ and $O(N^{-\Rea s})$ terms.
Since $g_{N+2}-g_{N+1}=O(N^{-\Rea t-1})$, the grouped boundary
tends to zero under \eqref{eq:ordered-domain}. The estimates are
locally uniform with an arbitrarily small exponent margin.
The remaining sum in \eqref{eq:finite-ordered} is
$4E_s(a)\zeta(t,b+1)+\Def(s,t;a,b)$.
Division by four proves the theorem.
\end{proof}

\subsection{What symmetrization removes}

The existing polarized determinant evaluates
$\tfrac12(\Ord(s,t;a,b)+\Ord(t,s;b,a))$.
The two one-index kernel differences satisfy the finite shift identity
\begin{equation}\label{eq:kernel-symmetry}
 \Def(s,t;a,b)+\Def(t,s;b,a)
 =-a^{-s}(b+1)^{-t}-(a+1)^{-s}b^{-t}.
\end{equation}
Indeed, shift the first positive kernel once, pair it with the
opposite negative kernel, and retain its missing zeroth term;
do the same with the other pair. This is first legitimate in an
absolute-convergence domain and then holds by continuation.
Combining \eqref{eq:kernel-symmetry} with \eqref{eq:ordered-main}
recovers the existing polarized formula. It also identifies exactly
which information the averaging has removed.

At a common shift and common order,
\begin{equation}\label{eq:diagonal-recovery}
 \Ord(s,s;a,a)=E_s(a)\zeta(s,a)
                  -\frac{a^{-2s}}2+
                     \frac{[a(a+1)]^{-s}}4.
\end{equation}
This is the known two-particle specialization of the incoming
single-zeta determinant \cite{repo:finite}. Thus the present theorem
extends that formula while preserving its normalization.

% End included source: sections/02_ordered
% Begin included source: sections/03_continuation
\section{Joint continuation, oblique poles, and directional values}
\label{sec:continuation}

\subsection{A convergent expansion in single Hurwitz zeta functions}

Set $\delta=a-b$ and $y=b+1$. For a sufficiently large integer $M$,
expand each shifted factor about $k+y$ on the tail $k\ge M$.
This gives
\begin{align}\label{eq:hurwitz-expansion}
 \Def(s,t;a,b)={}&
 \sum_{k=0}^{M-1}(k+y)^{-t}
       \bigl\{(k+a+2)^{-s}-(k+a)^{-s}\bigr\}\notag\\
 &+\sum_{j=1}^{\infty}\frac{(-1)^j(s)_j}{j!}
      \bigl\{(\delta+1)^j-(\delta-1)^j\bigr\}
           \zeta(s+t+j,M+y).
\end{align}
Choose $M$ so that both $|\delta+1|$ and $|\delta-1|$ are
strictly smaller than $|M+y|$; on any compact set a uniform strict
inequality is available. The factor $(s)_j/j!$ grows at most
polynomially in $j$ on compact $s$-sets, whereas the large-order
Hurwitz series supplies geometric decay after multiplication by
the displayed shift powers. The finitely many possible small-order
zeta poles are handled separately. Hence the tail is locally normally
convergent away from those poles, together with every fixed derivative.
This proves joint meromorphic continuation to all $(s,t)\in\C^2$.

\begin{theorem}[Complete polar data]\label{thm:poles}
The possible polar divisors of $\Ord$ are
\begin{equation}\label{eq:polar-divisors}
 t=1,\qquad s+t=1-j\quad(j=1,2,\ldots).
\end{equation}
Away from intersections, their residues in the indicated linear
coordinates are
\begin{align}\label{eq:polar-coefficients}
 \Res_{t=1}\Ord&=E_s(a),\notag\\
 \Res_{s+t=1-j}\Ord&=
 C_j(s;\delta):=\frac{(-1)^j(s)_j}{4j!}
              \bigl\{(\delta+1)^j-(\delta-1)^j\bigr\}.
\end{align}
These coefficient formulas continue across intersections. Every local
expression is a sum of terms with single linear denominators and a
holomorphic remainder; no product denominator is required. A listed
divisor can disappear on a special shift or order locus. The divisor
$s+t=1$ is absent identically.
\end{theorem}

\begin{proof}
In \eqref{eq:hurwitz-expansion} a pole occurs only when
$s+t+j=1$, with residue equal to its explicit coefficient, since
Hurwitz zeta has residue one. The $j=0$ binomial coefficient is zero
in the difference, so its would-be divisor cancels before
continuation. The other term in \eqref{eq:ordered-main} has only
$t=1$ as a possible pole. Local normal convergence of the remaining
tail proves the stated local decomposition.
\end{proof}

\begin{corollary}[Cancellation in the averaged sum]\label{cor:pole-cancellation}
Every oblique polar divisor cancels between $\Ord(s,t;a,b)$ and
$\Ord(t,s;b,a)$.
\end{corollary}
\begin{proof}
Equation \eqref{eq:kernel-symmetry} is entire in both orders. One
can also check residues directly: on $s+t=1-j$,
$(t)_j=(-1)^j(s)_j$, and replacing $\delta$ by $-\delta$
reverses the required coefficient. Thus the cancellation is an exact
functional identity, not a numerical relation between periods.
\end{proof}

\subsection{Four intersections at which every weight vanishes}

Both $w_n(0,a)$ and $w_n(-1,a)$ vanish for every $n$. The continued
ordered sum nevertheless has directional values at pairs of these
parameters. The origin gives the simplest explanation.

\begin{theorem}[Universal origin behavior]\label{thm:origin}
Near $(s,t)=(0,0)$,
\begin{equation}\label{eq:origin-germ}
 \Ord(s,t;a,b)=-\frac{s}{2(s+t)}+H(s,t;a,b),
 \qquad H\text{ holomorphic},\qquad H(0,0;a,b)=0.
\end{equation}
Consequently, if $\alpha+\beta\ne0$,
\begin{equation}\label{eq:origin-ray}
 \lim_{\varepsilon\to0}
 \Ord(\alpha\varepsilon,\beta\varepsilon;a,b)
        =-\frac{\alpha}{2(\alpha+\beta)}.
\end{equation}
The limit is independent of both spatial shifts.
\end{theorem}

\begin{proof}
The $j=1$ coefficient in \eqref{eq:polar-coefficients} is exactly
$-s/2$. All other divisors miss a sufficiently small neighborhood
of the origin. After subtracting this polar term, the remainder
is holomorphic. On the punctured slice $s=0$ the explicit
formula \eqref{eq:ordered-main} vanishes identically, so the
holomorphic remainder vanishes at the origin. Substitution of a
non-tangent line proves \eqref{eq:origin-ray}.
\end{proof}

The diagonal value $-1/4$ recovers the specialization anomaly already
recorded in \cite{repo:finite}. Formula \eqref{eq:origin-ray} is more
informative: the slice $s=0$ has limit zero, the slice $t=0$ has limit
$-1/2$, and general slopes interpolate rationally. These are limits
of the same meromorphic continuation, not different summation rules
chosen for the pointwise zero series.

\begin{theorem}[The complete $2\times2$ zero-weight table]\label{thm:zero-table}
Let $m,n\in\{0,1\}$ and $\alpha+\beta\ne0$. Along
$s=-m+\alpha\varepsilon$, $t=-n+\beta\varepsilon$,
\begin{equation}\label{eq:zero-ray}
 \lim_{\varepsilon\to0}\Ord(s,t;a,b)
       =\frac{\alpha}{\alpha+\beta}\,c_{m,n}(\delta),
\end{equation}
where
\begin{equation}\label{eq:zero-table}
 \begin{array}{c|cc}
 c_{m,n}(\delta)&n=0&n=1\\\hline
 m=0&-\tfrac12&\tfrac{\delta}{2}\\[2pt]
 m=1&-\tfrac{\delta}{2}&\tfrac{3\delta^2+1}{12}
 \end{array}.
\end{equation}
\end{theorem}
\begin{proof}
Only $j=m+n+1$ supplies a divisor through the given point.
Its coefficient $C_j(s;\delta)$ vanishes at $s=-m$ and has
derivative $c_{m,n}(\delta)$ there. The first three polynomials are
\[
 C_1=-\frac s2,\qquad
 C_2=\frac{\delta s(s+1)}2,\qquad
 C_3=-\frac{(3\delta^2+1)s(s+1)(s+2)}{12}.
\]
The restrictions $\Ord(0,t;a,b)$ and $\Ord(-1,t;a,b)$ are
identically zero meromorphic functions of $t$, directly from
\eqref{eq:ordered-main}. Thus the holomorphic remainder at the
chosen point vanishes. Expanding the one polar quotient proves
the table and the directional limit.
\end{proof}

No limit is asserted for directions with $\alpha+\beta=0$.
They lie on the polar hyperplane itself, where cancellation of a
numerator at one point does not define a restriction of the
meromorphic function to the whole line.

% End included source: sections/03_continuation
% Begin included source: sections/04_integer_gamma
\section{Finite polygamma evaluation and a digamma master}
\label{sec:integer}

The one-index reduction is useful only if its new kernel can be handled
explicitly. At integer orders it is finite; at confluence its apparent
singularities disappear; and the entire positive-order family is
generated by one function of two shifts.

\subsection{Partial fractions with all simple poles paired}

\begin{proposition}\label{prop:partial-fractions}
Let $p,q$ be positive integers and $\Rea x,\Rea y>0$. If $x\ne y$, put
$D=y-x$ and
\begin{align}
 A_k&=(-1)^{p-k}\binom{p+q-k-1}{q-1}D^{-(p+q-k)},\quad 1\le k\le p,
 \label{eq:pf-a}\\
 B_l&=(-1)^p\binom{p+q-l-1}{p-1}D^{-(p+q-l)},\quad 1\le l\le q.
 \label{eq:pf-b}
\end{align}
Then $A_1+B_1=0$ and
\begin{equation}\label{eq:finite-kernel}
 \Kern(p,q;x,y)=A_1\{\psi(y)-\psi(x)\}
 +\sum_{k=2}^{p}A_k\zeta(k,x)
 +\sum_{l=2}^{q}B_l\zeta(l,y).
\end{equation}
The empty sums are zero. At $x=y$ the value is $\zeta(p+q,x)$.
\end{proposition}

\begin{proof}
Expand $(X+D)^{-q}$ at $X=0$, and $X^{-p}$ at $X=-D$.
The respective principal parts of $X^{-p}(X+D)^{-q}$ have coefficients
\eqref{eq:pf-a} and \eqref{eq:pf-b}. Their difference from that rational
function is entire and vanishes at infinity, hence is zero.
Its $X^{-1}$ coefficient at infinity gives $A_1+B_1=0$.
Consequently the only two series which would diverge separately occur
as the convergent pair
\[
 A_1\sum_{n\ge0}\left(\frac1{n+x}-\frac1{n+y}\right)
 =A_1\{\psi(y)-\psi(x)\}.
\]
All higher powers sum absolutely. This proves \eqref{eq:finite-kernel}.
The confluent value follows directly from the original kernel.
\end{proof}

More generally, confluence comes with every Taylor coefficient:
\begin{equation}\label{eq:kernel-confluence}
 \Kern(s,t;x,x+h)=\sum_{r=0}^{\infty}
 \frac{(-1)^r(t)_r}{r!}h^r\zeta(s+t+r,x),\qquad |h|<|x|.
\end{equation}
Here $\Rea x>0$ and $\Rea(x+h)>0$; on compact sets a uniformly
smaller disk is understood.
The identity follows first for $\Rea(s+t)>1$ by the binomial theorem,
since $|n+x|\ge |x|$, and then meromorphically. It is preferable to
subtracting nearly cancelling negative powers of $D$ when computing
close shifts. The finite formula and \eqref{eq:classical-ladders}
establish a finite polygamma evaluation of every
$\Ord(p,q;a,b)$ with $p\ge1$, $q\ge2$.

\subsection{The exact parity left by equal shifts}

\begin{theorem}\label{thm:parity}
For $p\ge1$, $q\ge2$, $N=p+q$ and $\Rea a>0$, define
\[
 L_j=\begin{cases}\binom{N-j-1}{q-1},&j\le p,\\0,&j>p,\end{cases}
 \qquad
 R_j=\begin{cases}\binom{N-j-1}{p-1},&j\le q,\\0,&j>q.\end{cases}
\]
Then
\begin{align}
 \Ord(p,q;a,a)={}&E_p(a)\zeta(q,a+1)\notag\\
 &+\frac{(-1)^q}{4}\sum_{j=2}^{\max(p,q)}
 \{1-(-1)^{N-j}\}\{L_j+(-1)^jR_j\}\zeta(j,a+1)\notag\\
 &-\frac14\sum_{k=1}^{p}\binom{N-k-1}{q-1}
 \{(-1)^q(a+1)^{-k}+(-1)^{p-k}a^{-k}\}.
 \label{eq:ordered-parity}
\end{align}
Only orders $j$ with $p+q+j$ odd occur in the correction in the second
line. When $p=q$, that entire line vanishes.
\end{theorem}

\begin{proof}
Apply \eqref{eq:finite-kernel} to
$\Kern(p,q;a+2,a+1)-\Kern(p,q;a,a+1)$, so that $D=-1$ and $D=1$.
Use $\zeta(j,a)=a^{-j}+\zeta(j,a+1)$ and
$\zeta(j,a+2)=\zeta(j,a+1)-(a+1)^{-j}$, together with the same
shift recurrence for $\psi$. For $j\ge2$ the resulting coefficient is
$(-1)^q\{1-(-1)^{N-j}\}\{L_j+(-1)^jR_j\}/4$.
The simple-pole pairs contribute only the finite third line.
If $p=q$, then $L_j=R_j$: the second brace vanishes for odd $j$,
and the first vanishes for even $j$.
\end{proof}

This is a statement about explicit coefficients. It makes no claim
that the special values remaining in a particular formula are
arithmetically independent.

\begin{example}[The ordering survives]\label{ex:ordered-values}
Two equal-shift evaluations are
\begin{align}
 \Ord(2,3;1,1)&=\frac38\zeta(3)-\frac32\zeta(2)+\frac{33}{16},
 \label{eq:example-23}\\
 \Ord(3,2;1,1)&=\frac{31}{16}\zeta(2)-\frac{95}{32}.
 \label{eq:example-32}
\end{align}
Their average is the previously polarized example
$7\zeta(2)/32+3\zeta(3)/16-29/64$ in \cite{repo:finite}.
Confluent shifts give a different cancellation:
\begin{align}
 \Ord(2,2;1,2)&=\frac14\zeta(4)+\frac14\zeta(2)-\frac9{16},
 \label{eq:example-1122}\\
 \Ord(2,2;2,1)&=\frac7{36}\zeta(2)-\frac14\zeta(4)-\frac{11}{576}.
 \label{eq:example-2211}
\end{align}
The $\zeta(4)$ term is an exact confluence contribution erased when
the two assignments are added.
\end{example}

\subsection{One master function generates all positive orders}

Define the holomorphic divided difference
\[
 \Psi[x,y]=\begin{cases}
 \dfrac{\psi(y)-\psi(x)}{y-x},&x\ne y,\\[5pt]
 \psi'(x),&x=y.
 \end{cases}
\]
For $\Rea a,\Rea b>0$ set
\begin{equation}\label{eq:digamma-master}
 F(a,b)=-\frac{\psi(b+1)}{2a(a+1)}
 +\frac14\{\Psi[a+2,b+1]-\Psi[a,b+1]\}.
\end{equation}

\begin{theorem}[Digamma master]\label{thm:digamma-master}
For $p\ge1$ and $q\ge2$,
\begin{equation}\label{eq:master-derivatives}
 \Ord(p,q;a,b)=
 \frac{(-1)^{p+q-2}}{(p-1)!(q-1)!}
 \partial_a^{p-1}\partial_b^{q-1}F(a,b).
\end{equation}
For real $a,b>0$ and $|x|<a$, $|y|<b$ this is equivalently
\begin{equation}\label{eq:master-geometric}
 \sum_{p\ge1}\sum_{q\ge2}x^{p-1}y^{q-1}\Ord(p,q;a,b)
 =F(a-x,b-y)-F(a-x,b).
\end{equation}
Both apparent singularities $b=a\pm1$ in a quotient version of $F$
are removable.
\end{theorem}

\begin{proof}
The residue of $\Ord(1,t;a,b)$ at $t=1$ is
$E_1(a)=1/(2a(a+1))$. The constant of the Hurwitz zeta factor there is
$-\psi(b+1)$, while $\Kern(1,1;x,y)=\Psi[x,y]$ is convergent.
Thus \eqref{eq:digamma-master} is precisely
$\FP_{t=1}\Ord(1,t;a,b)$ with $a,b$ fixed.

Direct differentiation of the weights gives
\begin{equation}\label{eq:weight-derivatives}
 \partial_a^k w_n(s,a)=(-1)^k(s)_k w_n(s+k,a).
\end{equation}
For $q\ge2$,
\[
 \partial_b^{q-1}\Ord(1,t;a,b)
 =(-1)^{q-1}(t)_{q-1}\Ord(1,t+q-1;a,b).
\]
The $b$ derivative kills the residue, and the right side is regular at
$t=1$. Its value is $(-1)^{q-1}(q-1)!\Ord(1,q;a,b)$.
Differentiate in $a$ and use \eqref{eq:weight-derivatives} again.
This proves \eqref{eq:master-derivatives}; Taylor expansion proves
\eqref{eq:master-geometric}.

There is also a direct convergence interpretation. Summing the
geometric weight expansions produces $w_m(1,a-x)$ in the inner
index and $w_n(1,b-y)-w_n(1,b)$ in the outer index. The latter
difference has one additional power of decay. Their original ordered
sum is absolutely convergent on the stated polydisk. This proves that
\eqref{eq:master-geometric} is an ordinary generating identity.
\end{proof}

The slightly larger domain $a>0$, $b>-1$, $|x|<a$, $|y|<b+1$
also works: the outer summation starts at $n=1$. We retain the
smaller common shift domain elsewhere to keep the original definitions
uniform.

\subsection{Nonlinear harmonic weights and Gamma ratios}

Let $\Bell_r$ denote the complete exponential Bell polynomial, normalized by
\begin{equation}\label{eq:bell-definition}
 \exp\left(\sum_{j\ge1}x_j\frac{v^j}{j!}\right)
 =\sum_{r\ge0}\Bell_r(x_1,\ldots,x_r)\frac{v^r}{r!}.
\end{equation}
For $h_j(n,a)=\sum_{\ell=0}^{2}(n+a+\ell)^{-j}$, one has
\begin{equation}\label{eq:bell-weight}
 w_n(p,a)=\frac{\Gamma(n+a)}{\Gamma(n+a+3)}
 \frac{\Bell_{p-1}(h_1,1!h_2,\ldots,(p-2)!h_{p-1})}{(p-1)!}.
\end{equation}
For $p=1$ the Bell polynomial is $\Bell_0=1$.
Indeed $w_n(1,a)=\Gamma(n+a)/\Gamma(n+a+3)$, and the logarithmic
derivatives of this ratio are $(-1)^j(j-1)!h_j$.
Equation \eqref{eq:weight-derivatives} and the Bell formula for
successive derivatives prove \eqref{eq:bell-weight}.

This translates the ordered identities into explicit nonlinear
harmonic sums. For example, write
$P_n=(n+1)(n+2)(n+3)$ and
$h_{j,n}=H_{n+3}^{(j)}-H_n^{(j)}$. Then \eqref{eq:example-23} becomes
\begin{equation}\label{eq:harmonic-example}
 \sum_{0\le m<n}\frac{(n-m)^2h_{1,m}(h_{1,n}^2+h_{2,n})}{P_mP_n}
 =\frac34\zeta(3)-3\zeta(2)+\frac{33}{8}.
\end{equation}
All higher Bell combinations are supplied by the same finite master,
with no new summation coordinates. At integral shifts they involve
finite harmonic differences; at arbitrary shifts the $h_j$ are
finite polygamma differences by the usual recurrence.

\subsection{A Gamma primitive under common translation}

For $a,b>0$, $h\ge0$, define
$\Lambda_h(x)=\log\Gamma(x+h)-\log\Gamma(x)$.
Then
\begin{equation}\label{eq:gamma-primitive}
 \begin{split}
 \int_0^h\Def(1,1;a+c,b+c)\dd c
 ={}&\frac{\Lambda_h(b+1)-\Lambda_h(a+2)}{b-a-1}\\
 &-\frac{\Lambda_h(b+1)-\Lambda_h(a)}{b-a+1}.
 \end{split}
\end{equation}
At a zero denominator its quotient is the confluent derivative
$\Lambda_h'(x)=\psi(x+h)-\psi(x)$. To prove the identity, replace
each $\Kern(1,1)$ by its digamma divided difference; common
translation leaves its denominator constant, and the primitive of
$\psi$ is $\log\Gamma$. This also proves the holomorphic path version
in the right half-plane. For positive integral $h$, the Gamma ratios
are finite products, so the primitive becomes a finite logarithmic
identity. Further argument derivatives generate polygamma primitives
of all positive-order kernels.

% End included source: sections/04_integer_gamma
% Begin included source: sections/05_polylogarithms
\section{The cutoff polylogarithm and its disappearing depth}
\label{sec:polylog}

The ordered endpoint evaluation is accompanied by an explicit cutoff
identity. This is useful both as a proof of the endpoint cancellation
and as a finite transform for noninteger orders, one of the questions
raised in \cite{repo:finite}.

\subsection{An identity entire in the two orders}

For $|z|<1$ put
\begin{align}
 G_{s,t}(z)&=\sum_{0\le m<n}(n-m)^2
             w_m(s,1)w_n(t,1)z^{n+1},\label{eq:cutoff-g}\\
 B_{s,t}(z)&=\sum_{k\ge1}\frac{z^k}{k^s(k+1)^t},
 \qquad c_s=\frac{1-2^{-s}}2,\qquad \Theta=z\frac{\dd}{\dd z},
 \label{eq:cutoff-b}\\
 U_t(z)&=\frac{(1-z)^2}{z^2}\Li_t(z)+2-z^{-1}-2^{-t}.
 \label{eq:cutoff-u}
\end{align}
The last expression has a removable singularity at zero and equals
$2\sum_{n\ge0}w_n(t,1)z^{n+1}$. The same convention will apply to
the removable negative powers of $z$ in the next formula.

\begin{theorem}[Finite cutoff formula]\label{thm:cutoff}
For all $s,t\in\C$ and $|z|<1$,
\begin{align}
 G_{s,t}(z)={}&\frac{(1-z)^2}{2z^2}\Li_{t,s}(z,1)
 +\frac{z-1}{2z}\Li_{s+t}(z)
 +\frac14B_{t,s}(z)-\frac1{4z}B_{s,t}(z)\notag\\
 &+\frac12\left(c_s\Theta^2-\Theta-\frac12\right)U_t(z)
 +\frac12-2^{-t-2}.
 \label{eq:cutoff-master}
\end{align}
Every term is entire in $s,t$ on the open unit disk. At positive
integer orders the two $B$ functions are finite combinations of
ordinary polylogarithms. The only depth-two term in this representation
has the explicit factor $(1-z)^2$.
\end{theorem}

\begin{proof}
The inner finite identity of Section~\ref{sec:ordered}, with $a=1$,
becomes
\begin{equation}\label{eq:inner-harmonic}
 \sum_{m=0}^{n-1}(n-m)^2w_m(s,1)
 =H_{n+1}^{(s)}+c_s(n+1)^2-(n+1)-\frac12
   +\frac1{2(n+2)^s}.
\end{equation}
It also holds at $n=0$, since the right side then vanishes.
For $k\ge1$, the three generating identities needed for its harmonic
part are
\begin{align}
 \sum_{k\ge1}\frac{z^kH_k^{(s)}}{k^t}
   &=\Li_{t,s}(z,1)+\Li_{s+t}(z),\notag\\
 \sum_{k\ge1}\frac{z^kH_k^{(s)}}{(k+1)^t}
   &=z^{-1}\Li_{t,s}(z,1),\notag\\
 \sum_{k\ge1}\frac{z^kH_k^{(s)}}{(k+2)^t}
   &=z^{-2}\Li_{t,s}(z,1)-z^{-1}B_{s,t}(z).
 \label{eq:harmonic-generators}
\end{align}
They follow by writing the finite harmonic sum and changing the
outer index; the term subtracted in the last line is precisely the
last inner index in the shifted nested sum.

Multiply \eqref{eq:inner-harmonic} by $w_n(t,1)z^{n+1}$ and sum.
Its quadratic polynomial contributes
$\tfrac12(c_s\Theta^2-\Theta-\tfrac12)U_t$.
The harmonic contribution, by \eqref{eq:harmonic-generators}, is
\[
 \frac{(1-z)^2}{2z^2}\Li_{t,s}(z,1)
 +\frac12\Li_{s+t}(z)-\frac1{2z}B_{s,t}(z).
\]
The final term $1/(2(n+2)^s)$ contributes
\[
 \frac14B_{t,s}(z)-\frac1{2z}\Li_{s+t}(z)
 +\frac1{4z}B_{s,t}(z)+\frac12-2^{-t-2}.
\]
Adding gives \eqref{eq:cutoff-master}. On compact sets of orders,
every summand grows at most polynomially in its indices before the
exponential factor is applied. All rearrangements therefore converge
normally and prove the identity for every pair of complex orders.
\end{proof}

\subsection{Finite classical polylogarithms at integer orders}

For positive integers $p,q$, use the coefficients of
\eqref{eq:pf-a}--\eqref{eq:pf-b} with $D=1$. Then
\begin{equation}\label{eq:b-classical}
 B_{p,q}(z)=\sum_{k=1}^{p}A_k\Li_k(z)
 +\sum_{l=1}^{q}B_lz^{-1}\{\Li_l(z)-z\}.
\end{equation}
This follows by summing the rational partial fractions against $z^n$.
At $|z|<1$ the two simple-pole sums need not be paired for convergence,
although pairing is necessary before taking their individual endpoint
values. Since $\Theta\Li_\nu=\Li_{\nu-1}$, the differential term in
\eqref{eq:cutoff-master} is also finite in classical polylogarithms.
Thus the integer-order generator is explicitly in depth at most two,
and its depth-two coefficient is known exactly.

\begin{corollary}[Endpoint cancellation]\label{cor:cutoff-endpoint}
For positive integers $p\ge1$, $q\ge2$, the radial limit is
\begin{equation}\label{eq:cutoff-endpoint}
 \Ord(p,q;1,1)=c_p\zeta(q)-\frac12+2^{-p-2}
 +\frac14\{B_{q,p}(1)-B_{p,q}(1)\}.
\end{equation}
Together with \eqref{eq:b-classical}, interpreted through its paired
simple poles, this yields the parity formula
\eqref{eq:ordered-parity} at $a=1$.
\end{corollary}

\begin{proof}
Here $\Li_{q,p}(1,1)$ converges: the inner harmonic sum is bounded
for $p>1$ and grows only logarithmically for $p=1$, whereas $q\ge2$.
Its factor $(1-z)^2$ makes its contribution zero. The factor
$z-1$ similarly kills $\Li_{p+q}(z)$.
Finite summation by parts for the three moments of the single weight
gives
\[
 U_q(1)=1-2^{-q},\qquad
 (\Theta U_q)(1)=1,\qquad
 (\Theta^2U_q)(1)=2\zeta(q)-1.
\]
These moments converge for $q>1$. Substitution into
\eqref{eq:cutoff-master} gives \eqref{eq:cutoff-endpoint}.
Absolute convergence of the original double sum justifies its radial
limit as well.
\end{proof}

There is a broader analytic interpretation. If $\Rea t>1$ and
$\Rea(s+t)>0$, the original $G_{s,t}$ has the endpoint
$\Ord(s,t;1,1)$. The depth-two contribution still tends to zero
radially. Indeed
$\sum_{j<k}j^{-\Rea s}$ is
$O(1+k^{\max(1-\Rea s,0)}\log(k+1))$, and a direct comparison of
the resulting power series with a Laplace integral proves
$(1-z)^2\Li_{t,s}(z,1)\to0$ on this domain. Some of the remaining
separated terms can diverge and must be kept in the combination in
\eqref{eq:cutoff-master}. The identity of the collective endpoint
with the boundary formula is valid; replacing a divergent component
by an unannounced regularized value is unnecessary.

\subsection{A finite noninteger-order transform}

Writing $B_{s,t}$ as a one-index series already reduces depth, but a
standard special-function transform makes the answer more explicit.
Let
\[
 \Phi(z,v,c)=\sum_{k\ge0}\frac{z^k}{(k+c)^v},\qquad |z|<1,
\]
be the Lerch transcendent \cite{dlmf:lerch}. For $\Rea s,\Rea t>0$,
\begin{equation}\label{eq:b-beta-lerch}
 B_{s,t}(z)=\frac{z\Gamma(s+t)}{\Gamma(s)\Gamma(t)}
 \int_0^1u^{s-1}(1-u)^{t-1}
       \Phi(z,s+t,2-u)\dd u.
\end{equation}
To prove it, apply the beta identity
\begin{equation}\label{eq:feynman-basic}
 X^{-s}Y^{-t}=\frac{\Gamma(s+t)}{\Gamma(s)\Gamma(t)}
 \int_0^1\frac{u^{s-1}(1-u)^{t-1}}{[uX+(1-u)Y]^{s+t}}\dd u
\end{equation}
with $X=k$, $Y=k+1$ and sum $z^k$. The initial exchange is absolutely
convergent. Both sides then continue in their order variables from
that domain; the continued $B$ itself is entire for $|z|<1$.
Equations \eqref{eq:cutoff-master} and \eqref{eq:b-beta-lerch}
give a finite description of the cutoff generator for noninteger
orders using ordinary polylogarithms, one depth-two polylogarithm,
and two ordinary beta--Lerch transforms. No assumption that arbitrary
complex orders have finite integer partial fractions is made.

\subsection{Euler primitives}

For a function vanishing at zero, define
$\mathcal J f(z)=\int_0^z f(u)\dd u/u$ within the unit disk.
Termwise integration yields, for every integer $\ell\ge0$,
\begin{equation}\label{eq:euler-primitives}
 \mathcal J^\ell G_{s,t}(z)
 =\sum_{0\le m<n}\frac{(n-m)^2}{(n+1)^\ell}
             w_m(s,1)w_n(t,1)z^{n+1}.
\end{equation}
At positive integer orders these remain finite combinations of
depth-at-most-two polylogarithms and rational functions. One constructive
proof uses \eqref{eq:inner-harmonic}: the coefficient is a rational
function of $k=n+1$, times $H_k^{(p)}$, plus a rational function.
Dividing by $k^\ell$ introduces only additional integer-order poles
at $0,-1,-2$. Partial fractions reduce each term to one of
\eqref{eq:harmonic-generators} with shifted orders, or to its finite
index-shift analogues. Shifting $H_{k+c}^{(p)}$ changes it by a
finite rational sum, so no further nested index is introduced.
This proves closure under these primitives constructively; it does
not assert a minimal transcendental basis at the endpoint.

% End included source: sections/05_polylogarithms
% Begin included source: sections/06_stieltjes
\section{Ordinary Stieltjes identities and their primitives}
\label{sec:stieltjes}

Two complementary constructions produce every generalized Stieltjes
constant. The first subtracts a single asymptotic term from the original
ordered summand and retains ordinary convergence. The second takes
derivatives of a beta transform precisely on a hyperplane where the
two Hurwitz residues cancel.

\subsection{Subtracting before summing}

For $\Rea s>0$ and $\Rea a,\Rea b>0$, put
\[
 I_n(s,a)=\sum_{m=0}^{n-1}(n-m)^2w_m(s,a),\qquad
 c_2(t)=\frac{t(t+1)}2.
\]
Define the bracketed series
\begin{equation}\label{eq:subtracted-series}
 \mathcal U(s,t;a,b)=
 \sum_{n=1}^{\infty}
 \left[
 w_n(t,b)I_n(s,a)-E_s(a)c_2(t)(n+b)^{-t}
 \right].
\end{equation}
The subtraction is inside each summand.

\begin{theorem}[Convergent Stieltjes jets]\label{thm:subtracted-jets}
The series \eqref{eq:subtracted-series} is locally normally convergent
for $\Rea s>0$, $\Rea t>0$. On this domain,
\begin{equation}\label{eq:subtracted-identity}
 \mathcal U(s,t;a,b)=\frac14\Def(s,t;a,b)
       +(1-c_2(t))E_s(a)\zeta(t,b+1),
\end{equation}
with a removable singularity at $t=1$.
For $m\ge0$ define the ordinarily convergent logarithmic difference
\begin{equation}\label{eq:j-m}
 J_m(s;a,b)=\sum_{k=0}^{\infty}
 \frac{\{(k+a+2)^{-s}-(k+a)^{-s}\}
                 \log^m(k+b+1)}{k+b+1}.
\end{equation}
With $\mathcal U_m=\partial_t^m\mathcal U(s,1;a,b)$ and
$E=E_s(a)$, the exact identities are
\begin{align}
 \mathcal U_0&=\frac14J_0-\frac32E,\label{eq:u-zero}\\
 \mathcal U_1&=-\frac14J_1-\frac E2\{1+3\gamma_0(b+1)\},
 \label{eq:u-one}\\
 \mathcal U_m&=(-1)^m\left[
 \frac14J_m+\frac{3m}{2}E\gamma_{m-1}(b+1)
 -\frac{m(m-1)}2E\gamma_{m-2}(b+1)
 \right],\quad m\ge2.
 \label{eq:u-all}
\end{align}
Every $\mathcal U_m$ on the left is obtained by differentiating the
bracketed original series term by term.
\end{theorem}

\begin{proof}
The finite identity for $I_n$ gives
$I_n(s,a)=E_s(a)(n+b)^2+O(n)$ locally uniformly for $\Rea s>0$.
For the second difference,
\[
 w_n(t,b)=c_2(t)(n+b)^{-t-2}+O(n^{-\Rea t-3})
\]
uniformly on compact parameter sets. Their product after the
subtraction is $O(n^{-\Rea t-1})$; any fixed order derivative only
introduces powers of $\log n$, and a small compact-set margin
absorbs them. Thus the series is locally normal.

For $\Rea t>1$, subtract the Hurwitz series from the ordered
boundary identity. This proves \eqref{eq:subtracted-identity} there
and hence throughout the stated domain by analytic continuation.
At $t=1+v$,
\[
 1-c_2(1+v)=-\frac32v-\frac12v^2.
\]
Multiply this polynomial by \eqref{eq:stieltjes-convention}.
The $m$th derivative of $\Def(s,t;a,b)$ at $t=1$ is
$(-1)^mJ_m$, justified by its normally convergent combined series.
Extracting the constant, the linear coefficient, and the general
coefficient proves \eqref{eq:u-zero}--\eqref{eq:u-all}.
\end{proof}

These identities provide a uniform way of expressing Stieltjes
constants through logarithmically weighted ordered harmonic kernels.
They are not divergent double sums assigned values after the fact:
their defining brackets are summable before any continuation is used.

\begin{example}[An explicit nonlinear harmonic subtraction]
At $a=b=1$, $s=2$, the values are
\[
 E=\frac38,\qquad J_0=\frac{11}{4}-2\zeta(2),\qquad
 \mathcal U_0=\frac18-\frac12\zeta(2).
\]
For $s=3$ one instead obtains
$J_0=15/8-2\zeta(3)$ and
$\mathcal U_0=-3/16-\zeta(3)/2$.
In both cases the left side of \eqref{eq:subtracted-series} is
explicit in the harmonic Bell weights \eqref{eq:bell-weight}.
\end{example}

\subsection{The hard cutoff has a different constant}

Let $\Ord_N$ denote the original ordered sum with $n\le N$.
For $\Rea s>0$,
\begin{equation}\label{eq:hard-cutoff}
 \lim_{N\to\infty}
 \{\Ord_N(s,1;a,b)-E_s(a)\log(N+b)\}
 =\FP_{t=1}\Ord(s,t;a,b)-\frac32E_s(a),
\end{equation}
where the finite part keeps $s,a,b$ fixed. To prove it directly,
write the finite sum as
$\mathcal U_N(s,1;a,b)+E_s(a)\sum_{n=1}^N(n+b)^{-1}$.
The first term tends to \eqref{eq:u-zero}; the second, after
subtracting $E_s(a)\log(N+b)$, tends to $-E_s(a)\psi(b+1)$.
But
\begin{equation}\label{eq:fixed-order-fp}
 \FP_{t=1}\Ord(s,t;a,b)
 =-E_s(a)\psi(b+1)+\frac14J_0(s;a,b).
\end{equation}
This proves \eqref{eq:hard-cutoff}. Equivalently, the upper endpoints
in the exact finite summation by parts have limit $-3E_s(a)/2$.

For example,
\[
 F(1,1)=\frac{\gamma}{4}-\frac38,\qquad
 \lim_{N\to\infty}\left\{\Ord_N(1,1;1,1)-\frac14\log(N+1)\right\}
 =\frac{\gamma}{4}-\frac34.
\]
There is no contradiction: the two operations retain different
boundary constants. A diagonal Laurent coordinate $s=t$ would also
differentiate the residue $E_s(a)$ and is yet another declared
coordinate. These distinctions should be preserved in any integration
of the formulas into the manuscript.

\subsection{A beta transform on the cancelled polar hyperplane}

For $\Rea x,\Rea y>0$, $\Rea s,\Rea t>0$ and initially
$\Rea(s+t)>1$, summing
\eqref{eq:feynman-basic} gives
\begin{equation}\label{eq:k-beta}
 \Kern(s,t;x,y)=\frac{\Gamma(s+t)}{\Gamma(s)\Gamma(t)}
 \int_0^1u^{s-1}(1-u)^{t-1}
       \zeta(s+t,y+(x-y)u)\dd u.
\end{equation}
For $\Def$ take the difference of the two integrands before
continuing. The Hurwitz pole at $s+t=1$ then cancels identically.
This is a direct instance of the beta-parametrization method for
zeta products; compare \cite{shpot-paris}. The following all-order
Stieltjes identity uses that cancellation.

\begin{theorem}[Beta--Stieltjes difference identity]\label{thm:beta-stieltjes}
Suppose $0<\Rea\sigma<1$ and $\Rea a,\Rea b>0$. Set
\begin{align*}
 A_\pm(u)&=b+1+(a-b\pm1)u,\\
 T_m(\sigma;a,b)&=\sum_{k=0}^{\infty}
 \{(k+a+2)^{-\sigma}-(k+a)^{-\sigma}\}
 (k+b+1)^{\sigma-1}[-\log(k+b+1)]^m,
\end{align*}
and define
\begin{align}
 \ell_1(u)&=\log(1-u)+\psi(1)-\psi(1-\sigma),\notag\\
 \ell_j&=\psi^{(j-1)}(1)-\psi^{(j-1)}(1-\sigma),\qquad j\ge2.
 \label{eq:beta-ells}
\end{align}
Then, for every integer $m\ge0$,
\begin{equation}\label{eq:beta-stieltjes-all}
 \begin{split}
 T_m(\sigma;a,b)=\frac{\sin(\pi\sigma)}{\pi}
 \int_0^1&u^{\sigma-1}(1-u)^{-\sigma}\\[-2pt]
 {}\times&\sum_{j=0}^{m}\binom mj(-1)^j
 \Bell_{m-j}(\ell_1(u),\ldots,\ell_{m-j})
 \{\gamma_j(A_+(u))-\gamma_j(A_-(u))\}\dd u.
 \end{split}
\end{equation}
Both the sum and the integral converge ordinarily. The sum has
summands $O(k^{-2}\log^m k)$.
\end{theorem}

\begin{proof}
In the difference of \eqref{eq:k-beta}, put
$s=\sigma$, $t=1-\sigma+v$. Its integral is
\[
 \frac{\Gamma(1+v)}{\Gamma(\sigma)\Gamma(1-\sigma+v)}
 \int_0^1 u^{\sigma-1}(1-u)^{-\sigma+v}
 \{\zeta(1+v,A_+(u))-\zeta(1+v,A_-(u))\}\dd u.
\]
The expression in braces is holomorphic at $v=0$; its $j$th
derivative there is $(-1)^j(\gamma_j(A_+)-\gamma_j(A_-))$.
The factor outside the braces, divided by its value at $v=0$,
has logarithmic derivatives \eqref{eq:beta-ells}. The constant
factor is $\sin(\pi\sigma)/\pi$ by Gamma reflection.
The Bell product rule \eqref{eq:bell-definition} therefore gives
\eqref{eq:beta-stieltjes-all}.

All endpoint powers have real exponents greater than $-1$, uniformly
for $v$ in a sufficiently small disk. Additional logarithms are
integrable. The shifts $A_\pm(u)$ remain in a compact subset of the
right half-plane, so the differentiated Hurwitz difference is
bounded there. These facts justify differentiation under the
integral. On the series side, the shifted difference gains one
power of $k$, giving $O(k^{-2+\lvert\Rea v\rvert}\log^m k)$
uniformly near zero. Differentiating it gives exactly $T_m$.
\end{proof}

At $\sigma=\tfrac12$,
\begin{equation}\label{eq:half-beta-ells}
 \ell_1(u)=\log[4(1-u)],\qquad
 \ell_j=(-1)^j(j-1)!(2-2^j)\zeta(j),\quad j\ge2.
\end{equation}
Thus the transform has an arcsine weight and Bell coefficients
expressed entirely in ordinary zeta values. At $a=b=1$ there is
also the geometrically convergent single-zeta-derivative series
\begin{equation}\label{eq:beta-jet-zeta}
 T_m(\sigma;1,1)
 =-2\sum_{h=0}^{\infty}\frac{(\sigma)_{2h+1}}{(2h+1)!}
       \zeta^{(m)}(2h+2,2).
\end{equation}
Indeed the Hurwitz expansion of $\Def$ has only odd shift powers
when $a=b$, and along $s+t=1+v$ its zeta argument is $2h+2+v$.
The shift-to-tail ratio is $1/2$, so differentiating this expansion
is locally normal. In particular,
$T_0(\tfrac12;1,1)=-1/\sqrt2$; the next subsection explains its
elementary telescoping origin.

\subsection{Arcsine Gamma identities and Stieltjes primitives}

Some symmetric members simplify further and should be identified
as recurrence consequences, rather than presented as independent
priority claims. For $a>0$,
\begin{align}
 \int_0^1\frac{\psi(a+1+u)-\psi(a+1-u)}{\sqrt{u(1-u)}}\dd u
 &=\frac{\pi}{\sqrt{a(a+1)}},\label{eq:arcsine-psi}\\
 \int_0^1\frac{\log\Gamma(a+1+u)-\log\Gamma(a+1-u)}
                       {\sqrt{u(1-u)}}\dd u
 &=2\pi\log\frac{\sqrt a+\sqrt{a+1}}2.
 \label{eq:arcsine-gamma}
\end{align}
To prove these, replace $u$ by $1-u$ in the second member of each
difference inside its integral. The recurrences for $\psi$ and
$\Gamma$ then leave respectively $1/(a+u)$ and $\log(a+u)$.
The first arcsine integral is obtained by $u=(1+\cos\theta)/2$;
integrating it in $a$ and comparing as $a\to\infty$ gives the
second. This proof fixes the normalization of the logarithm,
including the factor $2$ in its denominator.

From $\gamma_m(a+1)-\gamma_m(a)=-\log^m(a)/a$ one likewise obtains
the all-order identity
\begin{equation}\label{eq:arcsine-stieltjes}
 \int_0^1\frac{\gamma_m(a+1+u)-\gamma_m(a+1-u)}
                      {\sqrt{u(1-u)}}\dd u
 =-\int_0^1\frac{\log^m(a+u)}{(a+u)\sqrt{u(1-u)}}\dd u.
\end{equation}
For instance, its $m=1$ value is
\begin{equation}\label{eq:arcsine-gamma-one}
 -\frac{\pi}{\sqrt{a(a+1)}}
 \log\frac{4a(a+1)}{(\sqrt a+\sqrt{a+1})^2}.
\end{equation}
For a direct evaluation, set $A=a+\tfrac12$, $B=\tfrac12$,
$D=\sqrt{A^2-B^2}$ and $r=B/(A+D)$. Then
$A+B\cos\theta=(A+D)(1+2r\cos\theta+r^2)/2$.
The Fourier series of its logarithm and the Poisson kernel of its
reciprocal give
\[
 \int_0^\pi\frac{\log(A+B\cos\theta)}{A+B\cos\theta}\dd\theta
 =\frac{\pi}{D}\log\frac{2D^2}{A+D}.
\]
Their series converge absolutely since $0<r<1$; substitution is
\eqref{eq:arcsine-gamma-one} with the required negative sign.

There is also a canonical primitive without an arbitrary periodic
ambiguity. Define
$S_m(a)=\int_1^a\gamma_m(x)\dd x$ for $a>0$.
The elementary integral
\[
 \int_1^2\zeta(1+v,a)\dd a
 =\frac{\zeta(v,1)-\zeta(v,2)}{v}=\frac1v
\]
shows, coefficient by coefficient, that $\int_1^2\gamma_m(a)\dd a=0$.
Integrating the shift recurrence consequently gives
\begin{equation}\label{eq:stieltjes-primitive-recurrence}
 S_m(a+1)-S_m(a)=-\frac{\log^{m+1}a}{m+1}.
\end{equation}
The same symmetry argument yields
\begin{equation}\label{eq:arcsine-primitive}
 \int_0^1\frac{S_m(a+1+u)-S_m(a+1-u)}{\sqrt{u(1-u)}}\dd u
 =-\frac1{m+1}\int_0^1\frac{\log^{m+1}(a+u)}{\sqrt{u(1-u)}}\dd u.
\end{equation}
For $m=0$, $S_0(a)=-\log\Gamma(a)$, and this reduces exactly to
\eqref{eq:arcsine-gamma}. Equations
\eqref{eq:beta-stieltjes-all} and \eqref{eq:arcsine-primitive}
exhibit two different sources of all-order identities: cancellation
of a zeta pole in a beta transform, and integration of the elementary
shift recurrence.

% End included source: sections/06_stieltjes
% Begin included source: sections/07_localization
\section{All-rank localization of ordered finite-difference sums}
\label{sec:localization}

The incoming finite-difference report proves determinant identities for
symmetrized assignments of orders and shifts, and explicitly asks which
coordinates are needed for an individual ordered assignment
\cite[Sections 3 and 9]{repo:finite}.  We give a finite construction for
arbitrary rank.  Its central fact is a support cancellation: after the
differences are transferred to the ordered Vandermonde polynomial, two
isolated free labels annihilate the coefficient.  The construction retains
the individual assignment of every order and shift.

\subsection{The ordered sum and its exact transported coefficient}

Fix an integer $r\ge2$ and put $d=2r-2$.  For
$\mathbf s=(s_1,\ldots,s_r)\in\C^r$ and
$\Rea a_i>0$, define
\begin{align}
 w_i(n)&=\frac1{d!}\sum_{j=0}^{d}(-1)^j\binom dj
                       (n+a_i+j)^{-s_i},
                       \label{loc:weights}\\
 \mathcal Z_r(\mathbf s;\mathbf a)
   &=\sum_{0\le n_1<\cdots<n_r}
          \Delta(\mathbf n)^2\prod_{i=1}^r w_i(n_i),
                       \label{loc:ordered-sum}
\end{align}
where $\Delta(\mathbf x)=\prod_{i<v}(x_v-x_i)$.
All powers use the principal logarithm in the right half-plane.  We shall
distinguish two domains:
\[
 \Omega=\{\Rea s_i>1,\ \Rea a_i>0\text{ for every }i\},
 \qquad
 \Omega_0=\{\Rea s_i>d+1,\ \Rea a_i>0\text{ for every }i\}.
\]

\begin{theorem}[Finite ordered reduction in every rank]
\label{loc:main}
The series \eqref{loc:ordered-sum} converges absolutely and locally normally
on $\Omega$.  On $\Omega_0$ it is an explicit finite sum of expressions
\begin{equation}
 \sum_{0\le N_1<\cdots<N_q}
    Q(N_1,\ldots,N_q)
    \prod_{\ell=1}^q\prod_{i\in B_\ell}
                      (N_\ell+A_i)^{-s_i},
 \qquad \Rea A_i>0,
 \label{loc:reduced-form}
\end{equation}
multiplied, where necessary, by finite products of shifted powers.
Here $Q$ is a polynomial, the label sets $B_\ell$ are consecutive blocks,
and
\[
 q\le\left\lfloor\frac{r+1}{2}\right\rfloor.
\]
The contribution anchored at the lower boundary has at most
$\lfloor r/2\rfloor$ free indices.  Equations
\eqref{loc:sector-polynomial}, \eqref{loc:free-sector}, and
\eqref{loc:anchored-sector} below give all coefficients and shifts by
finite sums.  Every expression in this construction converges absolutely
on $\Omega_0$.  Their collective sum has a holomorphic continuation to
$\Omega$, equal there to \eqref{loc:ordered-sum}.
\end{theorem}

The number $q$ counts summation variables in the displayed representation.
The theorem makes no assertion that this number is minimal, or that any
special values obtained from different sectors are arithmetically
independent.

For the convergence statement, the finite-difference integral is
\begin{equation}
 w_i(n)=\frac{(s_i)_d}{d!}
       \int_{[0,1]^d}
       (n+a_i+u_1+\cdots+u_d)^{-s_i-d}\dd\mathbf u.
 \label{loc:difference-integral}
\end{equation}
The sign is positive because $d$ is even.  On compact subsets of the
parameter domain this gives
$|w_i(n)|\le C(1+n)^{-\Rea s_i-d}$.  For nonnegative integer nodes,
\begin{equation}
 |\Delta(\mathbf n)|^2\le\prod_{i=1}^r(1+n_i)^d.
 \label{loc:vandermonde-bound}
\end{equation}
The resulting product majorant proves normal convergence on $\Omega$.
If the differences in \eqref{loc:weights} are expanded before summation,
each factor is only bounded by $C(1+n_i)^{-\Rea s_i}$; the same argument
then proves absolute convergence on $\Omega_0$.  This is the initial
domain for the coefficient transport.

On integer tuples define the chamber polynomial
\begin{equation}
 P(\mathbf x)=\ind_{\{0\le x_1<\cdots<x_r\}}
                                      \Delta(\mathbf x)^2
 \label{loc:chamber}
\end{equation}
and the finite stencil
\begin{equation}
 C_r(\mathbf k)=
 \sum_{\mathbf j\in\{0,\ldots,d\}^r}
       (-1)^{j_1+\cdots+j_r}\prod_{i=1}^r\binom d{j_i}
       P(\mathbf k-\mathbf j).
 \label{loc:transport-coefficient}
\end{equation}
Expanding \eqref{loc:weights}, setting $k_i=n_i+j_i$, and collecting the
common shifted powers gives the exact identity
\begin{equation}
 \mathcal Z_r(\mathbf s;\mathbf a)
 =\frac1{(d!)^r}\sum_{k_1,\ldots,k_r\ge0}
               C_r(\mathbf k)\prod_{i=1}^r(k_i+a_i)^{-s_i}
 \qquad\text{on }\Omega_0.
 \label{loc:transport}
\end{equation}
Expanded absolute convergence justifies every regrouping.  There is no
unaccounted upper endpoint: the regrouping is performed on the full
absolutely convergent series.  Also, $C_r(\mathbf k)=0$ if any $k_i<0$,
since every nonzero term in \eqref{loc:transport-coefficient} has
$k_i=x_i+j_i\ge0$.

\subsection{Collision blocks and the two-singleton cancellation}

The shifted tuple $\mathbf k$ need not be increasing.  If
$C_r(\mathbf k)\ne0$, however, some stencil term has a strictly increasing
tuple $\mathbf k-\mathbf j$.  Necessarily
\begin{equation}
 k_{i+1}-k_i\ge1-d\qquad(1\le i<r).
 \label{loc:gap-lower}
\end{equation}
Partition the labels $1,\ldots,r$ into maximal consecutive blocks by
placing a cut precisely when
\begin{equation}
 k_{i+1}-k_i>d.
 \label{loc:cut}
\end{equation}
The internal adjacent gaps therefore belong to the finite set
$\{1-d,\ldots,d\}$.  Across a cut, every stencil assignment remains
strictly ordered, because
$k_{i+1}-j_{i+1}-k_i+j_i\ge d+1-d=1$.

The first block is called \emph{anchored} when $k_1<d$; every other block
is free.  If $k_1\ge d$, all blocks are free.  Only the first block can
be anchored.  Indeed, the chamber condition in \eqref{loc:chamber}
requires just $x_1\ge0$ together with strict order, which already implies
nonnegativity of the later coordinates.

\begin{lemma}[Two free singleton blocks give zero]
\label{loc:singleton}
If the block decomposition of $\mathbf k$ contains two free blocks of
size one, then $C_r(\mathbf k)=0$.
\end{lemma}

\begin{proof}
Let $i$ be one free singleton, and fix every stencil variable except
$j_i$.  All chamber indicators are independent of $j_i$: its neighboring
cuts remain ordered for every shift, and it cannot change the lower
boundary condition.  As a polynomial in $x_i=k_i-j_i$, the Vandermonde
square has degree $d=2r-2$, with leading coefficient
\[
 \Delta(\mathbf x_{\widehat i})^2
       =\prod_{\substack{p<v\\p,v\ne i}}(x_v-x_p)^2.
\]
Consequently
\[
 \sum_{j_i=0}^d(-1)^{j_i}\binom d{j_i}\Delta(\mathbf x)^2
       =d!\,\Delta(\mathbf x_{\widehat i})^2.
\]
Let $v$ be the second free singleton.  Its indicators are still
independent of $j_v$, but the remaining polynomial has degree
$2(r-2)=d-2$ in $x_v$.  Its $d$th backward difference is zero.  Summing
over the other stencil variables preserves this cancellation.
\end{proof}

At most one surviving free block can therefore be a singleton.  If
there are $q$ free blocks and no anchor, they use at least $2q-1$ labels;
hence $q\le\lfloor(r+1)/2\rfloor$.  If an anchor is present, it uses at
least one additional label, giving $r\ge2q$ and
$q\le\lfloor r/2\rfloor$.  This proves the support part of
Theorem~\ref{loc:main}.  We now turn it into a finite formula without
leaving any block-coordinate choices implicit.

\subsection{An explicit finite construction of every sector}

Choose a consecutive block composition
$\mathcal B=(B_1,\ldots,B_p)$ of $\{1,\ldots,r\}$, and write
$b_\ell=|B_\ell|$ and $\ell(i)$ for the block containing $i$.  In each
block choose integer offsets $u_i$ with first offset zero and
\begin{equation}
 u_{i+1}-u_i\in\{1-d,\ldots,d\}
                    \quad(i,i+1\in B_\ell).
 \label{loc:offsets}
\end{equation}
There are finitely many choices.  Put
\[
 k_i=t_{\ell(i)}+u_i,
 \qquad c_\ell=d+u_{\max B_\ell},
\]
and introduce the finite internal indicator
\[
 \varepsilon_{\mathcal B,\mathbf u}(\mathbf j)
 =\prod_{\ell=1}^p\prod_{\substack{i,i+1\in B_\ell}}
       \ind_{\{u_{i+1}-u_i+j_i-j_{i+1}\ge1\}}.
\]
Discard a pattern if this indicator is zero for every stencil tuple.
For every retained pattern, an admissible tuple gives
\begin{equation}
 u_i\ge i-\min B_\ell-d\quad(i\in B_\ell),
 \qquad c_\ell\ge b_\ell-1\ge0.
 \label{loc:offset-feasibility}
\end{equation}
To see the first inequality, sum the internal strict inequalities from
the first label of the block to $i$, then use $j_i-j_{\min B_\ell}\ge-d$.
The second is its instance at the last label.

Define the following polynomial in the block bases:
\begin{align}
 \mathcal C_{\mathcal B,\mathbf u}(\mathbf t)
 ={}&\sum_{\mathbf j\in\{0,\ldots,d\}^r}
       (-1)^{|\mathbf j|}\prod_{i=1}^r\binom d{j_i}
       \varepsilon_{\mathcal B,\mathbf u}(\mathbf j)
       \prod_{i<v}
       \bigl(t_{\ell(v)}-t_{\ell(i)}+u_v-u_i-j_v+j_i\bigr)^2.
 \label{loc:sector-polynomial}
\end{align}
It has integer coefficients and satisfies
\begin{equation}
 \deg_{t_\ell}\mathcal C_{\mathcal B,\mathbf u}
                         \le 2b_\ell(r-b_\ell).
 \label{loc:sector-degree}
\end{equation}
Pairs of labels in the same block do not depend on its base; each of
the $b_\ell(r-b_\ell)$ cross-block pairs contributes at most degree two.

For an unanchored sector the base conditions are exactly
\[
 t_1\ge d,\qquad t_{\ell+1}>t_\ell+c_\ell.
\]
Set
\begin{equation}
 e_1=d,\qquad e_\ell=d+\sum_{v<\ell}c_v,
 \qquad t_\ell=N_\ell+e_\ell.
 \label{loc:free-bases}
\end{equation}
The conditions become $0\le N_1<\cdots<N_p$.  Moreover,
$e_\ell+u_i\ge0$ by \eqref{loc:offset-feasibility}, so every shifted
argument stays in the right half-plane.  This sector contributes exactly
\begin{equation}
 \frac1{(d!)^r}
 \sum_{0\le N_1<\cdots<N_p}
       \mathcal C_{\mathcal B,\mathbf u}(\mathbf N+\mathbf e)
       \prod_{\ell=1}^p\prod_{i\in B_\ell}
                    (N_\ell+e_\ell+u_i+a_i)^{-s_i}.
 \label{loc:free-sector}
\end{equation}
The stencil polynomial is exactly $C_r(\mathbf k)$ on this sector:
cross-block order and the lower-bound condition are automatic, leaving
precisely the indicators in \eqref{loc:sector-polynomial}.

For an anchored sector choose the finite integer
$h=k_1\in\{0,\ldots,d-1\}$ and set $t_1=h$.  Let
$\mathcal C_{\mathcal B,\mathbf u;h}(t_2,\ldots,t_p)$ be
\eqref{loc:sector-polynomial} with $t_1=h$ and the additional indicator
$\ind_{\{j_1\le h\}}$ inserted in the summand.  Discard patterns with
no admissible stencil tuple under this extra constraint.  If $p\ge2$,
set
\begin{equation}
 e_2=h+c_1+1,\qquad
 e_\ell=e_2+\sum_{v=2}^{\ell-1}c_v\quad(\ell\ge3),
 \qquad t_\ell=N_\ell+e_\ell.
 \label{loc:anchored-bases}
\end{equation}
The free variables satisfy $0\le N_2<\cdots<N_p$.  Feasibility of the
anchored block gives $h+u_i\ge0$ for its labels, and in particular
$e_2=d+h+u_{\max B_1}+1\ge d+1$.  Together with
\eqref{loc:offset-feasibility}, this also ensures positive real parts for
all the free shifted arguments.  The anchored sector is
\begin{align}
 &\frac1{(d!)^r}
   \prod_{i\in B_1}(h+u_i+a_i)^{-s_i}
   \sum_{0\le N_2<\cdots<N_p}
     \mathcal C_{\mathcal B,\mathbf u;h}
                      (N_2+e_2,\ldots,N_p+e_p)
 \notag\\[-2pt]
 &\hspace{44mm}\times
       \prod_{\ell=2}^p\prod_{i\in B_\ell}
                   (N_\ell+e_\ell+u_i+a_i)^{-s_i}.
 \label{loc:anchored-sector}
\end{align}
For $p=1$, the sum and the free product are absent, so the contribution
is a finite scalar expression.

Every possible nonzero coefficient in \eqref{loc:transport} belongs to
exactly one of these sectors: its cuts determine the composition, the
first coordinate in each block determines the base, and subtraction
determines the offsets.  Conversely, the indicated inequalities give
exactly those cuts and internal gaps.  There are finitely many
compositions, offsets, and anchored values $h$.  Summing
\eqref{loc:free-sector} and \eqref{loc:anchored-sector} over them, and
discarding the zero sectors of Lemma~\ref{loc:singleton}, is therefore
the promised finite construction.

For completeness, each separated sector converges absolutely on
$\Omega_0$.  If a free block has $b$ labels, its polynomial degree in
the associated variable is at most $2b(r-b)$, while the sum of the real
parts of its denominator exponents is greater than $b(d+1)$.  Their
difference satisfies
\[
 b(d+1)-2b(r-b)=b(2b-1)\ge1,
\]
with strict improvement supplied by $\Rea s_i>d+1$.  A product
majorant over unrestricted nonnegative free variables is summable.
The same bounds apply locally uniformly to the parameters.  This proves
the convergence assertions and completes the proof of
Theorem~\ref{loc:main}.

If a representation without polynomial numerators is preferred, choose
one denominator $(N+A)^{-s}$ in each free block.  Expanding
$N^k=((N+A)-A)^k$ removes any monomial numerator by finitely many integer
decrements of that exponent.  The preceding degree estimate proves
absolute convergence of these terms on $\Omega_0$ as well, and no new
summation variable is introduced.

\begin{remark}[Collective continuation and separated convergence]
\label{loc:continuation-caution}
The original series is holomorphic on $\Omega$.  Since the finite
localized expression agrees with it on the nonempty subdomain
$\Omega_0$, the former gives a holomorphic continuation of the latter
collectively.  The degree estimate above does not prove that each
separated sector, or each individual exponent-shifted term, converges
on all of $\Omega$.  Finite cancellation between offset patterns may be
needed there.  This distinction is part of the theorem's domain
statement, not an appeal to formal rearrangement outside convergence.
\end{remark}

\subsection{Rank three: the complete two-index contact part}

For $r=3$, $d=4$, write $f_i(k)=(k+a_i)^{-s_i}$.  The unanchored
composition $(1,2)\mid(3)$ is described by
\[
 k_1\ge4,\qquad h=k_2-k_1\in\{-3,\ldots,4\},
 \qquad k_3-k_2>4.
\]
Summing the isolated third stencil first, exactly as in
Lemma~\ref{loc:singleton}, gives
\begin{align}
 C_3(k_1,k_2,k_3)&=24A(h),\label{loc:rank3-contact}\\
 A(h)&=\sum_{j,k=0}^4(-1)^{j+k}\binom4j\binom4k
                  (h+j-k)^2\ind_{\{h+j-k>0\}}.
                  \label{loc:rank3-A}
\end{align}
This entire integer contact table has a short generating expression:
\begin{equation}
 \sum_{h\in\Z}A(h)z^h
   =(1-z^{-1})^4(1-z)^4\frac{z(1+z)}{(1-z)^3}
   =z^{-3}(1+z)(1-z)^5.
 \label{loc:rank3-generating}
\end{equation}
Indeed, $z(1+z)/(1-z)^3=\sum_{h\ge1}h^2z^h$, and the two finite
stencils multiply this series by the displayed Laurent factors.
One may first work in $0<|z|<1$, where every intermediate series
converges.  The final Laurent polynomial proves
\[
\begin{array}{c|rrrrrrr}
 h&-3&-2&-1&0&1&2&3\\\hline
 A(h)&1&-4&5&0&-5&4&-1
\end{array}
\qquad A(h)=0\quad(|h|\ge4).
\]

The resulting two-index contribution of $(1,2)\mid(3)$ is
\begin{equation}
 \mathcal T_{12\mid3}
  =\frac1{24^2}\sum_{h=-3}^3 A(h)
       \sum_{0\le n<m}
       f_1(n+4)f_2(n+4+h)f_3(m+8+h).
 \label{loc:rank3-first}
\end{equation}
The other unanchored two-block composition, $(1)\mid(2,3)$, gives
\begin{equation}
 \mathcal T_{1\mid23}
  =\frac1{24^2}\sum_{h=-3}^3 A(h)
       \sum_{0\le n<m}
       f_1(n+4)f_2(m+8)f_3(m+8+h).
 \label{loc:rank3-second}
\end{equation}
Every remaining surviving sector has at most one free index.  Thus
\begin{equation}
 \mathcal Z_3=\mathcal T_{12\mid3}+\mathcal T_{1\mid23}
                         +\mathcal E_3,
 \label{loc:rank3-decomposition}
\end{equation}
where $\mathcal E_3$ is exactly the finite sum prescribed by
\eqref{loc:free-sector}--\eqref{loc:anchored-sector} over the remaining
compositions and anchor choices.  More explicitly, it consists of the
unanchored one-block sectors, all anchored one-block sectors, and the
anchored two-block sectors; anchored three-block sectors and
unanchored three-block sectors vanish by Lemma~\ref{loc:singleton}.
This is a finite formula for its one-index correction, with no
undetermined coefficient.

The contact contribution is not vacuous: the admissible tuple
$(5,6,20)$ has $h=1$ and
\begin{equation}
 C_3(5,6,20)=24A(1)=-120.
 \label{loc:rank3-nonzero}
\end{equation}
This proves the existence of two-index support in the construction.
Further identities can still simplify the complete sum for special
choices of the orders or shifts; \eqref{loc:rank3-nonzero} is not a
minimal-depth assertion.

\subsection{General integer stencils and stronger annihilation}

The support mechanism depends on polynomial moments of the stencil,
rather than on consecutive binomial coefficients.  The next result
retains $d=2r-2$, even when the stencil has a larger support.

\begin{corollary}[Finite integer stencils]
\label{loc:general-stencil}
Let $L\ge d$ be an integer and, for each label $i$, choose coefficients
$\alpha_{i,\lambda}\in\C$ supported on $0\le\lambda\le L$ such that
\begin{equation}
 \sum_{\lambda=0}^L\alpha_{i,\lambda}\lambda^k=0
                      \qquad(0\le k<d).
 \label{loc:stencil-moments}
\end{equation}
Replace \eqref{loc:weights} by
\[
 w_{i,\alpha}(n)=\sum_{\lambda=0}^L
                     \alpha_{i,\lambda}(n+a_i+\lambda)^{-s_i}.
\]
The corresponding ordered sum has the conclusions of
Theorem~\ref{loc:main}, with the same bounds
$\lfloor(r+1)/2\rfloor$ and $\lfloor r/2\rfloor$ on its free indices.
If every stencil also annihilates degree $d$, the bounds improve to
$\lfloor r/2\rfloor$ for the unanchored part and
$\lfloor(r-1)/2\rfloor$ for the anchored part.
\end{corollary}

\begin{proof}
The exact transported coefficient is now
\begin{equation}
 C_\alpha(\mathbf k)=
   \sum_{\boldsymbol\lambda\in\{0,\ldots,L\}^r}
          \left(\prod_i\alpha_{i,\lambda_i}\right)
          P(\mathbf k-\boldsymbol\lambda).
 \label{loc:general-coefficient}
\end{equation}
The factor $(d!)^{-r}$ is absent unless it was included in the chosen
coefficients.  Expanded absolute convergence on $\Omega_0$ justifies
the same transport.  A nonzero coefficient has adjacent gaps at least
$1-L$.  Make a cut at gaps greater than $L$, and anchor the first
block when $k_1<L$.

At a free singleton $i$, a degree-$d$ polynomial is reduced to its
leading coefficient times
\[
 \mu_i=\sum_{\lambda=0}^L
                 \alpha_{i,\lambda}(-\lambda)^d.
\]
This follows by expanding each power of $k_i-\lambda$ and using
\eqref{loc:stencil-moments}.  A second free singleton then occurs in
a polynomial of degree $d-2$, so its stencil annihilates the term.
The previous counting argument gives the original bounds.  If each
$\mu_i=0$, even the first free singleton annihilates the coefficient.
Every free block then uses at least two labels; with or without an
anchor this gives exactly the two improved bounds.

The full finite sector construction is unchanged after replacing the
geometric thresholds $d$ by $L$, the shift range by $\{0,\ldots,L\}$,
and the binomial stencil product by $\prod_i\alpha_{i,\lambda_i}$.
In particular internal offset increments belong to
$\{1-L,\ldots,L\}$, $c_\ell=L+u_{\max B_\ell}\ge b_\ell-1$, and
the unanchored first base is $L$.  The anchored first base has finitely
many choices $0\le h<L$, with the condition $\lambda_1\le h$.
Feasibility gives $u_i\ge i-\min B_\ell-L$, so the same base
substitutions keep all shifted arguments in the right half-plane.
The polynomial degree bound \eqref{loc:sector-degree} does not change.
Consequently $\Omega_0$, with $d=2r-2$, still suffices for convergence
of every separated sector, independently of $L$.

Finally Taylor's formula with a uniform remainder, applied to
$(n+a_i+\lambda)^{-s_i}$ through degree $d-1$ in $\lambda$, gives
\[
 w_{i,\alpha}(n)=O\bigl((1+n)^{-\Rea s_i-d}\bigr)
\]
locally uniformly in the parameters.  The polynomial terms vanish by
\eqref{loc:stencil-moments}; the finite support bounds the remainders.
Equation~\eqref{loc:vandermonde-bound} proves original normal
convergence on $\Omega$, and the collective continuation follows as
before.
\end{proof}

Every consecutive finite difference of order strictly greater than
$2r-2$ satisfies the stronger hypothesis.  For nonintegral rational
stencils, transport onto one integer lattice would instead require
separate residue classes; global moment cancellation alone does not
justify the integer-sector argument above.

The transport itself is discrete summation by parts.  Discrete residues
provide a classical framework for rational summability
\cite{chen-singer}.  Here the specific additional ingredient is the
support of the ordered chamber coefficient and the loss of two
Vandermonde degrees after a singleton is removed.  That mechanism works
with arbitrary complex orders and unequal shifts on the stated genuine
convergence domains, and resolves the finite-coordinate representation
question for this family in \cite{repo:finite}.

% End included source: sections/07_localization
% Begin included source: sections/08_gamma_basis
\section{A Gamma integral basis and eventual reconstruction at every fixed gap}
\label{gb:section}

The preceding sections develop identities by localizing a finite-difference
operator at its boundary.  We now apply a different analytic construction
to a remaining reconstruction problem in the repository: the optimal
initial moments of two unequal even harmonic tails.  The inherited
fixed-gap reduction expresses the obstruction as a rational determinant
of size half the gap \cite{repo-gamma}.  We construct its entire null
space by ordinary Gamma integrals, determine the exact normalization,
and obtain a nonzero determinant asymptotic.  Consequently the optimal
prefix holds eventually on \emph{every fixed-gap ray}.  This differs from
the previously established eventual result with the smaller order fixed
\cite{repo:quartic}:
here both orders tend to infinity while their difference remains fixed.

The Gamma recurrence, Stirling expansion, and Fourier-transform setting
are classical \cite{dlmf:gamma,dlmf-stirling,ramanujan-transform}.
The contributions established below are the explicit basis for this
particular null equation, its exact bridge to the inherited rational
determinant, the fixed-gap asymptotic, and its harmonic-number correction.
This is a statement of the advance relative to the inspected material;
it does not assert historical priority.  Throughout the section, rank
means rank of the specified rational \emph{formal coordinate matrix}.
No arithmetic independence of evaluated zeta values or moments is
assumed or concluded.

\subsection{The inherited moments and the determinant to be studied}

For integers $1\le a<b$, write
\[
 p=2a,\qquad q=2b,\qquad m=b-a\ge1,
 \qquad x_n=n+\tfrac12.
\]
Thus $b$ is half the larger order and $m$ is half the gap.  Define
\begin{equation}\label{gb:tails}
 A_p(n)=\zeta(p,n+1)=\zeta(p)-H_n^{(p)},\qquad
 D_{p,q}(s)=\sum_{n=0}^{\infty}A_p(n)A_q(n)x_n^{-s}.
\end{equation}
The defining series converges absolutely for $\Rea s>3-p-q$.
Let $D_r=D_{p,q}(-2r)$, using meromorphic continuation when necessary.
These points are regular.  Indeed the centered Hurwitz expansion has
the form
\[
 A_p(n)\sim\frac{x_n^{1-p}}{p-1}
 +\sum_{k\ge1}\frac{B_{2k}(1/2)}{(2k)!}
       (p)_{2k-1}x_n^{1-p-2k}.
\]
Subtracting finitely many terms from the product continues
\eqref{gb:tails} successively to larger half-planes.  Its possible poles
lie at $s=3-p-q-2\ell$, $\ell\ge0$, all odd integers.  In particular
the even arguments used in $D_r$ are regular.

The inherited two-tail calculus supplies blocks $C_j$ and the mutually
inverse Bernoulli transforms
\begin{align}
 D_r&=\sum_{j=0}^r\frac{\binom{2r}{2j}}{2j+1}
          B_{2r-2j}(1/2)C_j,\label{gb:forward}\\
 C_j&=\sum_{r=0}^j\binom{2j+1}{2r}2^{-2(j-r)}D_r.
 \label{gb:inverse}
\end{align}
Here $B_n(x)$ denotes a Bernoulli polynomial and $B_n=B_n(0)$, with
$B_1=-1/2$.  For $j\ge b$ the late coordinates are
\begin{equation}\label{gb:late}
 C_j=\sum_{k=0}^{b-1}a_{jk}Z_k,\qquad
 Z_k=\zeta(2k),\qquad Z_0=-\tfrac12.
\end{equation}
Their exact rational coefficients, also inherited, are as follows.
Put $d=2j+1$ and $e=d-p-q+2k+1$.  If $e<0$ set $a_{jk}=0$;
otherwise set
\begin{equation}\label{gb:source-entry}
 a_{jk}=B_e\left\{
 \frac{\binom{d-q+1}{p-2k}}{d-q+1}
 +\frac{\binom{d-p+1}{q-2k}}{d-p+1}\right\}.
\end{equation}
Each binomial summand is omitted unless its lower index lies between
$2$ and its upper index.  In particular, a lower index of zero is
\emph{omitted}.  This convention is essential: it determines the
constant subtraction in the differential equation below.  Equations
\eqref{gb:forward}--\eqref{gb:source-entry} are inputs from
\cite[Section~2, inherited moment coordinates]{repo-gamma}.

Define the $b\times b$ late-block matrix
\begin{equation}\label{gb:M}
 M_{a,b}=(a_{b+r,k})_{0\le r,k<b}.
\end{equation}
On the ordered formal labels
\[
 C_0,\ldots,C_{b-1},\quad Z_0,\ldots,Z_{b-1},
\]
the coefficient matrix of $C_0,\ldots,C_{2b-1}$ is
$\operatorname{diag}(I_b,M_{a,b})$.  The forward transform
\eqref{gb:forward} is triangular with diagonal $1/(2r+1)$.
Consequently the determinant of the complete matrix for the first
$q=2b$ moments $D_0,\ldots,D_{2b-1}$ is
\begin{equation}\label{gb:full-prefix}
             \frac{\det M_{a,b}}{(4b-1)!!}.
\end{equation}
Thus the optimal-prefix question is exactly the nonvanishing of
\eqref{gb:M}.  If it is nonzero, those first $2b$ moments generate
every later moment in the inherited formal coordinate system, and
evaluation gives the corresponding rational generation statements for
the actual constants.  Fewer than $2b$ rows cannot have rank $2b$ on
these $2b$ specified labels.

\subsection{The fixed-gap recurrence and its exact normalization}

For the remainder of the section fix $m\ge1$.  To distinguish the
recurrence coefficients from both the tails and Bernoulli numbers, put
\begin{equation}\label{gb:recurrence-factors}
 \mathfrak a_r=\begin{cases}
 0,&1\le r\le m,\\
 (2r)^{\underline{2m}},&r>m,
 \end{cases}
 \qquad
 \mathfrak b_N=(2N-1)^{\underline{2m}},\qquad
 d_N=2(\mathfrak a_N+\mathfrak b_N)\quad(N>m).
\end{equation}
The underline denotes a falling factorial.  All $d_N$ in this range
are strictly positive.  Let
\begin{equation}\label{gb:beta}
 \beta_k=\frac{2|B_{2k}|}{(2k)!}
        =\frac{4\zeta(2k)}{(2\pi)^{2k}}>0,
 \qquad k\ge1.
\end{equation}
The rational fundamental array is defined by
\begin{align}
 C^{[m]}_{r,j}&=\delta_{rj}&& (1\le r,j\le m),\label{gb:initial}\\
 C^{[m]}_{N,j}&=\frac1{d_N}\sum_{r=1}^{N-1}
   \beta_{N-r}(\mathfrak a_r+\mathfrak b_N)C^{[m]}_{r,j}
       &&(N>m,\ 1\le j\le m).\label{gb:recurrence}
\end{align}
The superscript separates this array from the moment blocks $C_j$.
For $b>m$ write
\begin{equation}\label{gb:F}
                  F_{b,m}=(C^{[m]}_{b+i,j})_{1\le i,j\le m}.
\end{equation}

For an even formal series $R(z)=\sum_{n\ge1}r_nz^{2n}$, define
\begin{equation}\label{gb:cot-operator}
 \mathcal T_mR
 =\bigl(R^{(2m)}-R^{(2m)}(0)\bigr)\cot(z/2)
                      +\bigl(R\cot(z/2)\bigr)^{(2m)}.
\end{equation}
Both terms combine to an odd formal Taylor series: the subtraction
removes the possible pole in the first product, while $R\cot(z/2)$
is an odd Taylor series.  Expanding
\[
 \cot(z/2)=\frac2z-\sum_{k\ge1}\beta_kz^{2k-1}
\]
shows that, for $N>m$, its coefficient at $z^{2N-2m-1}$ is
\begin{equation}\label{gb:coefficient}
 d_Nr_N-\sum_{r=1}^{N-1}
      \beta_{N-r}(\mathfrak a_r+\mathfrak b_N)r_r.
\end{equation}
The value $\mathfrak a_m=0$ records precisely the subtracted
$R^{(2m)}(0)$.  Since $d_N>0$, the first $m$ coefficients are
arbitrary and determine all later ones uniquely.  The formal null
space of $\mathcal T_m$ has dimension exactly $m$.

\begin{proposition}[Inherited boundary reduction, with normalization]
\label{gb:bridge}
For $b>m\ge1$,
\begin{equation}\label{gb:bridge-formula}
 \boxed{\displaystyle
 \det M_{b-m,b}
 =\frac{(-1)^{b(b-1)/2}}{4^bb!}
       \left(\prod_{N=m+1}^{b+m}d_N\right)\det F_{b,m}.}
\end{equation}
This is the complementary-minor theorem of
\cite[Theorem~\texttt{gap:complementary}]{repo-gamma}.
\end{proposition}

\begin{proof}
We reproduce the normalization argument so that the subsequent
Gamma determinant can be checked against the original coordinates.
Set $h(t)=\coth(t/2)$ and
\[
 \mathcal L_mQ=(Q^{(2m)}-Q^{(2m)}(0))h+(Qh)^{(2m)}.
\]
Let $K_{m,b}$ map the coefficients $(q_1,\ldots,q_b)$ of
$Q(t)=\sum_{n=1}^bq_nt^{2n}$ to the coefficients of
$t,t^3,\ldots,t^{2b-1}$ in $\mathcal L_mQ$.
Write $h(t)=\sum_{k\ge0}h_kt^{2k-1}$, where $h_0=2$ and
$h_k=2B_{2k}/(2k)!$ for $k\ge1$.  In row $r$, indexed from zero,
put $N=r+m+1$.  The entry in column $n$ is
\[
 h_{N-n}(\mathfrak a_n+\mathfrak b_N)\quad(n\le N),
\]
and is zero for $n>N$.  In \eqref{gb:source-entry}, substitution
$j=b+r$, $k=b-n$ gives Bernoulli index $2(N-n)$ and the upper/lower
binomial pairs
\[
 (2r+2,\,2n-2m),\qquad(2r+2m+2,\,2n).
\]
Multiplication by $2(2n)!/(2r+1)!$ therefore gives exactly the
displayed operator entry, including the omission when $n\le m$.
If $M^{\mathrm{rev}}$ denotes column reversal, then
\[
 K_{m,b}=
 \operatorname{diag}_{0\le r<b}\!\left(\frac2{(2r+1)!}\right)
 M_{b-m,b}^{\mathrm{rev}}
 \operatorname{diag}_{1\le n\le b}\!\bigl((2n)!\bigr),
\]
so
\begin{equation}\label{gb:operator-conjugacy}
 \det K_{m,b}=(-1)^{b(b-1)/2}4^bb!\det M_{b-m,b}.
\end{equation}

The relation $h_k=(-1)^{k-1}\beta_k$ for $k\ge1$ and
\eqref{gb:recurrence} imply that
\[
 Q_j(t)=\sum_{n\ge1}(-1)^{n-j}C^{[m]}_{n,j}t^{2n},
 \qquad1\le j\le m,
\]
are formal null solutions of $\mathcal L_m$, with initial
$m\times m$ coefficient block equal to the identity.  Replace
the first $m$ columns of $K_{m,b}$ by the corresponding linear
combinations given by the truncations of these solutions through
$n=b$.  This changes no determinant.  The first $b-m$ rows of
those columns vanish.  Their bottom $m$ rows equal $-TQ_+$, where
$T$ is lower triangular with diagonal $d_{b+1},\ldots,d_{b+m}$ and
\[
 (Q_+)_{ij}=(-1)^{b+i-j}C^{[m]}_{b+i,j}.
\]
Indeed each truncated null equation is the negative of its omitted
future terms, and its newest future coefficient has factor $d_{b+i}$.
Moving the first $m$ columns to the end leaves the other diagonal
block triangular with diagonal $d_{m+1},\ldots,d_b$.  It follows that
\[
 \det K_{m,b}=(-1)^{m(b-m)+m}
       \left(\prod_{N=m+1}^{b+m}d_N\right)\det Q_+.
\]
Factoring its row and column signs gives
$\det Q_+=(-1)^{mb}\det F_{b,m}$.  The combined exponent
$2mb-m(m-1)$ is even.  Equation \eqref{gb:operator-conjugacy}
now proves \eqref{gb:bridge-formula}.
\end{proof}

\subsection{An explicit Gamma integral basis}

Define the positive even weight
\begin{equation}\label{gb:W}
 W_m(x)=\left|\Gamma\!\left(\tfrac14+\tfrac{ix}{2}\right)\right|^{4m}
 \qquad(x\in\mathbb R)
\end{equation}
and the odd multipliers
\begin{equation}\label{gb:P}
 P_j(x)=\sinh\!\bigl(\pi(m-j)x\bigr)\quad(1\le j<m),
 \qquad P_m(x)=\tanh(\pi x).
\end{equation}
Their order is chosen to agree with their distinct decay scales.

\begin{theorem}[Gamma basis and exact top moments]\label{gb:basis}
For $1\le j\le m$, the integral
\begin{equation}\label{gb:R}
 R_j(z)=4\sin(z/2)\int_0^\infty
                  \sinh(xz)P_j(x)W_m(x)\dd x
\end{equation}
defines an even holomorphic function in $|\Rea z|<j\pi$.
Its germ at zero satisfies $\mathcal T_mR_j=0$ and
\begin{equation}\label{gb:top-derivative}
 R_j^{(2m)}(0)=0\quad(j<m),\qquad
 R_m^{(2m)}(0)=(-1)^{m+1}2^{2m+2}\pi^m.
\end{equation}
The $m$ germs form a basis of the formal null space of
\eqref{gb:cot-operator}.  Their Taylor radii are exactly $j\pi$.
\end{theorem}

The contour argument below proves the null equation and the exact
derivatives.  The basis assertion will follow from the positive
initial determinant in \cref{gb:normalization}; the exact Taylor
radii will follow independently from \cref{gb:coefficient-asymptotic}.
This order of proof avoids assuming either linear independence or
the location of a singularity from the integral representation alone.

\begin{proof}[Contour proof]
Stirling's formula gives, as $x\to+\infty$,
\[
 W_m(x)=2^m(2\pi)^{2m}x^{-m}e^{-m\pi x}(1+O(x^{-2})).
\]
Together with \eqref{gb:P}, this proves absolute convergence in
the stated strips, locally uniform convergence of all differentiated
integrals, and evenness of \eqref{gb:R}.  We use the usual meromorphic
Gamma function and the integer-power continuation
\[
 W_m(y)=\left[
 \Gamma(\tfrac14+iy/2)\Gamma(\tfrac14-iy/2)\right]^{2m}.
\]
There is no branch choice for this continuation; the absolute-value
notation in \eqref{gb:W} is used only on the real line.

For any one of the allowed odd multipliers $P$, put
\[
 G_P(t)=\int_{\mathbb R}e^{ixt}P(x)W_m(x)\dd x,
 \qquad Q_P(t)=2\sinh(t/2)G_P(t).
\]
These integrals and all their derivatives converge absolutely for
real $t$ and define analytic germs.  Since $G_P$ is odd, $Q_P$ is
even and vanishes at zero.  With $D=d/dt$, let
\[
 E_P(t)=e^t(D+\tfrac12)^{2m}G_P(t)
                 -e^{-t}(D-\tfrac12)^{2m}G_P(t).
\]
The identities $Q_P=(e^{t/2}-e^{-t/2})G_P$ and
$Q_P\coth(t/2)=(e^{t/2}+e^{-t/2})G_P$ yield, by differentiation,
\begin{equation}\label{gb:contour-algebra}
 \coth(t/2)Q_P^{(2m)}(t)
       +(Q_P(t)\coth(t/2))^{(2m)}
                    =\frac{E_P(t)}{\sinh(t/2)}.
\end{equation}

In the first integral defining $E_P$, substitute $y=x-i$;
its contour is $\mathbb R-i$.  In the second substitute $y=x+i$;
its contour is $\mathbb R+i$.  The factors $e^{\pm t}$ disappear
under these substitutions.  Gamma recurrence gives exactly
\begin{equation}\label{gb:gamma-cancellation}
 (iy-\tfrac12)^{2m}W_m(y+i)
 =(iy+\tfrac12)^{2m}W_m(y-i)=2^{2m}V_m(y),
\end{equation}
where
\[
 V_m(y)=\left[
 \Gamma(\tfrac34+iy/2)\Gamma(\tfrac34-iy/2)\right]^{2m}.
\]
For example the first equality shifts
$\Gamma(-1/4+iy/2)$ to $\Gamma(3/4+iy/2)$; its recurrence
factor is $iy/2-1/4$.  The second equality has the opposite sign
before taking the even power, hence the same factor $2^{2m}$.

Every allowed multiplier satisfies $P(y+i)=P(y-i)$.  Both shifted
integrands are therefore the same meromorphic function
\[
 \mathcal F_P(y)=2^{2m}e^{iyt}V_m(y)P(y+i).
\]
The nearest Gamma poles of $V_m$ are $y=\pm3i/2$, outside
$|\Im y|\le1$.  For $P(x)=\sinh(\pi kx)$, $1\le k<m$,
the multiplier is $P(y+i)=(-1)^k\sinh(\pi ky)$, an entire
function.  Uniform Stirling estimates on the closed strip give
a polynomial times $e^{-\pi(m-k)|\Rea y|}$ as a bound for the
integrand, with a constant depending on the fixed real $t$.
For the tanh multiplier, $P(y+i)=\tanh(\pi y)$; away from its
finite poles it is bounded at the two ends of the strip, and
the decay is a polynomial times $e^{-m\pi|\Rea y|}$.
Thus the vertical sides of the rectangle between the two horizontal
contours vanish as their real parts tend to infinity.  Its bottom
side runs from left to right, so bottom minus top equals
$+2\pi i$ times the enclosed residue sum.

In the sinh case there are no enclosed poles and $E_P(t)=0$.
In the tanh case the only enclosed poles are the simple poles
$y=\pm i/2$.  Each tanh residue is $1/\pi$, and
$V_m(i/2)=V_m(-i/2)=\pi^m$.  Hence
\[
 \Res_{y=\pm i/2}\mathcal F_P(y)
      =2^{2m}\pi^{m-1}e^{\mp t/2},
\]
and
\begin{equation}\label{gb:residue-result}
 E_{\tanh}(t)=i\,2^{2m+2}\pi^m\cosh(t/2).
\end{equation}
All shifts have now been justified using pole-free Gamma factors,
the two explicitly listed tanh poles, and decaying vertical sides.

The second summand on the left of \eqref{gb:contour-algebra} is
analytic at zero.  Its first summand has Laurent coefficient
$2Q_P^{(2m)}(0)$ at $t^{-1}$.  Comparing this coefficient shows
that $Q_P^{(2m)}(0)=0$ for a sinh multiplier and
$Q_P^{(2m)}(0)=i2^{2m+2}\pi^m$ for tanh.  Equations
\eqref{gb:contour-algebra} and \eqref{gb:residue-result} therefore
give the exact germ identity
\[
 (Q_P^{(2m)}-Q_P^{(2m)}(0))\coth(t/2)
                  +(Q_P\coth(t/2))^{(2m)}=0.
\]
Finally set $R_P(z)=iQ_P(iz)$.  Pairing positive and negative
$x$ in $G_P(iz)$ gives precisely \eqref{gb:R}.  Rotation uses
$\coth(iz/2)=-i\cot(z/2)$ and yields
$\mathcal T_mR_P=0$.  Since
$R_P^{(2m)}(0)=i(-1)^mQ_P^{(2m)}(0)$, it also gives
\eqref{gb:top-derivative}.
\end{proof}

\begin{corollary}[Exact ordinary Gamma moments]\label{gb:moment-identities}
Put $\phi_n(x)=\Im(x+i/2)^{2n}$.  The Taylor coefficients of
\eqref{gb:R} are
\begin{equation}\label{gb:t}
 t_{n,j}:=[z^{2n}]R_j(z)
      =\frac4{(2n)!}\int_0^\infty\phi_n(x)P_j(x)W_m(x)\dd x.
\end{equation}
For every $m\ge1$,
\begin{align}
 \int_0^\infty\phi_m(x)\sinh(\pi kx)W_m(x)\dd x&=0
                  &&(1\le k<m),\label{gb:orthogonality}\\
 \int_0^\infty\phi_m(x)\tanh(\pi x)W_m(x)\dd x
          &=(-1)^{m+1}2^{2m}\pi^m.\label{gb:evaluation}
\end{align}
All these integrals converge absolutely.
\end{corollary}

\begin{proof}
The elementary generating identity
\[
 \sin(z/2)\sinh(xz)=
       \sum_{n\ge1}\frac{\phi_n(x)}{(2n)!}z^{2n}
\]
and local uniform exponential domination justify coefficient
extraction.  Now apply \eqref{gb:top-derivative}.
\end{proof}

The vanishing moments have a concrete polynomial interpretation.
Since $\phi_m(x)/x$ has degree $2m-2$ and its positive zeros are
\[
                  \frac12\cot\frac{r\pi}{2m},\qquad1\le r<m,
\]
equation \eqref{gb:orthogonality} is an explicit family of
multiple-orthogonality relations against positive weights on
$(0,\infty)$.  The zeros follow by writing
$x+i/2=|x+i/2|e^{i\theta}$, with $0<\theta<\pi/2$, and solving
$\sin(2m\theta)=0$.  These relations do not imply that every
later consecutive coefficient determinant has one fixed sign.

\subsection{A positive initial determinant and the exact Gamma representation}

Write
\[
 T_m=(t_{n,j})_{1\le n,j\le m},\qquad
 \Delta(y_1,\ldots,y_m)=\prod_{r<s}(y_s-y_r).
\]

\begin{proposition}[Positive normalization]\label{gb:normalization}
The initial determinant is
\begin{equation}\label{gb:detT}
                 \det T_m=(-1)^{m(m-1)/2}K_m,
\end{equation}
where $K_m>0$ has the absolutely convergent integral representation
\begin{equation}\label{gb:Km}
\begin{split}
 K_m={}&\frac{2^{(m-1)(m-2)/2}}{m!}
             \left(\prod_{n=1}^m\frac{4n}{(2n)!}\right)\\
 &\times\int_{(0,\infty)^m}
       \Delta(x_1^2,\ldots,x_m^2)
       \Delta(\cosh\pi x_1,\ldots,\cosh\pi x_m)\\
 &\hspace{38mm}\times
       \prod_{r=1}^m x_rW_m(x_r)\tanh(\pi x_r)
                    \dd x_1\cdots\dd x_m.
\end{split}
\end{equation}
Consequently the $R_j$ form a basis of the full formal null space,
and for every $b>m$,
\begin{equation}\label{gb:exact-Gamma-determinant}
 \boxed{\displaystyle
 \det F_{b,m}=
  \frac{(-1)^{m(m-1)/2}}{K_m}
        \det(t_{b+i,j})_{1\le i,j\le m}.}
\end{equation}
\end{proposition}

\begin{proof}
The polynomial $\phi_n(x)$ has leading term $nx^{2n-1}$ and
only lower odd powers thereafter.  A triangular row transformation
therefore gives
\[
 \det T_m=\left(\prod_{n=1}^m\frac{4n}{(2n)!}\right)
       \det\left(\int_0^\infty x^{2n-1}P_j(x)W_m(x)\dd x\right)_{n,j=1}^m.
\]
The continuous Cauchy--Binet identity writes the remaining determinant
as $1/m!$ times the integral of the two alternants
$\det(x_r^{2n-1})_{n,r}$ and $\det(P_j(x_r))_{r,j}$ against
$\prod_rW_m(x_r)\dd x_r$.  This identity follows directly by
expanding the two determinants, applying Fubini, and permuting the
integration variables; absolute convergence justifies each operation.
The first alternant is $\prod_rx_r\,\Delta(x_1^2,\ldots,x_m^2)$.

For the second alternant put $u=\pi x$, $v=\cosh u$, and use
\[
              \frac{\sinh(ku)}{\tanh u}=vU_{k-1}(v),
\]
where $U_{k-1}$ is the Chebyshev polynomial of the second kind,
with degree $k-1$ and leading coefficient $2^{k-1}$.
After factoring $\tanh(\pi x_r)$ from row $r$, the column
degrees are $m-1,m-2,\ldots,1,0$.  Their leading-coefficient
product is $2^{(m-1)(m-2)/2}$.  Reversing the columns thus gives
\[
\begin{split}
 \det(P_j(x_r))={}&(-1)^{m(m-1)/2}2^{(m-1)(m-2)/2}\\
 &\times\prod_r\tanh(\pi x_r)
              \Delta(\cosh\pi x_1,\ldots,\cosh\pi x_m).
\end{split}
\]
This proves \eqref{gb:detT}--\eqref{gb:Km}.  For distinct
positive $x_r$, the two Vandermondes have the same sign, since
both $x\mapsto x^2$ and $x\mapsto\cosh(\pi x)$ are strictly
increasing.  The integrand in \eqref{gb:Km} is strictly positive
away from the collision hyperplanes.  Its exponential growth from
the cosh Vandermonde is at most
$\prod_re^{(m-1)\pi x_r}$, while each $W_m(x_r)$ supplies
$e^{-m\pi x_r}$ times a polynomial factor.  At zero all factors
are bounded or vanish.  This proves convergence and $K_m>0$.

All $R_j$ are null germs, and their first $m$ coefficient columns
are independent by \eqref{gb:detT}.  The formal null space has
dimension $m$ by \eqref{gb:coefficient}, so they form a basis.
Their coefficient vectors obey \eqref{gb:recurrence}; consequently
\[
                  (t_{b+i,j})_{i,j=1}^m=F_{b,m}T_m.
\]
Taking determinants proves \eqref{gb:exact-Gamma-determinant}.
\end{proof}

\subsection{A Laplace-moment lemma with controlled error}

The asymptotic argument must retain one term beyond the leading
Laplace phase.  At every even mode that leading phase is real and
its imaginary part vanishes.  The following lemma supplies the
necessary next term without requiring an analytic continuation of
the error in a real density.

\begin{lemma}\label{gb:Laplace}
Let $m\ge1$ be fixed, $R>0$, and $a\in\mathbb R\setminus\{0\}$.
Suppose a real, locally bounded function $f$ on $[0,\infty)$ satisfies
\[
             f(x)=ax^{-m}e^{-Rx}(1+O(x^{-2}))\qquad(x\to\infty).
\]
As integers $N\to\infty$,
\begin{equation}\label{gb:Laplace-formula}
 \int_0^\infty(x+i/2)^Nf(x)\dd x
 =ae^{iR/2}\frac{\Gamma(N-m+1)}{R^{N-m+1}}
       \left(1+\frac{imR}{2(N-m)}+O(N^{-2})\right).
\end{equation}
The implicit constant may depend on $f,m,R$.
\end{lemma}

\begin{proof}
Choose fixed $X>1$ beyond which the asserted density estimate holds.
The integral over $[0,X]$ is $O(C^N)$ for a fixed $C$, hence
negligible compared with
$S_N=\Gamma(N-m+1)R^{-N+m-1}$ by Stirling's formula.
The absolute error made on $[X,\infty)$ by replacing $f(x)$ with
$ax^{-m}e^{-Rx}$ is bounded by a constant times
\[
             \int_X^\infty|x+i/2|^Nx^{-m-2}e^{-Rx}\dd x.
\]
Split this integral at $x=N^{2/3}$.  On the lower part it is at most
$\exp(\tfrac23N\log N+O(N))$, whereas
$S_N=\exp(N\log N+O(N))$.  This part is smaller than any fixed
inverse power of $N$ times $S_N$.  On the upper part,
\[
 \frac{|x+i/2|^N}{x^N}
   =\left(1+\frac1{4x^2}\right)^{N/2}
   \le\exp\!\left(\frac{N}{8x^2}\right)
\]
is uniformly bounded.  Its integral is therefore at most a constant
times $\Gamma(N-m-1)R^{-N+m+1}=O(N^{-2}S_N)$.
This estimate uses only real-variable bounds on the density error.

For the model integral set $y=x+i/2$.  After factoring $e^{iR/2}$,
its integrand is
$y^{N-m}(1-i/(2y))^{-m}e^{-Ry}$ on the ray
$[X,\infty)+i/2$.  Shift that ray to $[X,\infty)$ inside
$\Rea y\ge X$.  No pole is crossed.  The short vertical segment
at $X$ contributes $O(C^N)$, and the vertical segment at infinity
vanishes for each $N$ because of $e^{-R\Rea y}$.
Uniformly for real $y\ge X$,
\[
                 (1-i/(2y))^{-m}=1+\frac{im}{2y}+O(y^{-2}).
\]
Integration of the first two terms, extending their lower limit to
zero at another $O(C^N)$ cost, gives
\[
 ae^{iR/2}\left\{
 \frac{\Gamma(N-m+1)}{R^{N-m+1}}
 +\frac{im}{2}\frac{\Gamma(N-m)}{R^{N-m}}\right\}.
\]
The last term has absolute size $O(N^{-2}S_N)$ by the same Gamma
ratio estimate.  This proves \eqref{gb:Laplace-formula}.
\end{proof}

Define the constants
\begin{equation}\label{gb:mode-constants}
 \lambda_j=j\pi,\qquad \xi_j=(j\pi)^{-2},\qquad
 A_m^*=2^m(2\pi)^{2m},\qquad
 a_j=\begin{cases}A_m^*/2,&j<m,\\A_m^*,&j=m.
 \end{cases}
\end{equation}
Conjugate Gamma asymptotics contain only even inverse powers after
the leading factor.  Thus
\begin{equation}\label{gb:density-asymptotic}
 P_j(x)W_m(x)=a_jx^{-m}e^{-\lambda_jx}(1+O(x^{-2})).
\end{equation}
The further exponential terms in sinh or tanh are absorbed into the
indicated error.  In particular, there is no $x^{-1}$ correction.

\begin{proposition}[Distinct coefficient scales and exact radii]
\label{gb:coefficient-asymptotic}
Put $d_j=m+\ind_{2\mid j}$ and
\begin{equation}\label{gb:Lj}
 L_j=\begin{cases}
 4a_j\sin(j\pi/2)\lambda_j^{m-1},&j\text{ odd},\\
 2ma_j\cos(j\pi/2)\lambda_j^m,&j\text{ even}.
 \end{cases}
\end{equation}
All $L_j$ are nonzero, and
\begin{equation}\label{gb:t-asymptotic}
 t_{n,j}=L_j\lambda_j^{-2n}(2n)^{-d_j}(1+O(n^{-1})).
\end{equation}
The Taylor radius of $R_j$ is exactly $\lambda_j=j\pi$.
\end{proposition}

\begin{proof}
Apply \cref{gb:Laplace} to \eqref{gb:density-asymptotic}, take
the imaginary part, and multiply by $4/(2n)!$.  If $j$ is odd,
$e^{i\lambda_j/2}=\pm i$ and the leading term survives.  Since
$\Gamma(2n-m+1)/\Gamma(2n+1)\sim(2n)^{-m}$, this gives the
first line of \eqref{gb:Lj}.  If $j$ is even, the leading phase
is $\pm1$ and its imaginary part is zero.  The next term in
\eqref{gb:Laplace-formula} gives instead
\[
 \frac{4a_j\cos(j\pi/2)}{(2n)!}
 \frac{m}{2}\frac{\Gamma(2n-m)}{\lambda_j^{2n-m}},
\]
which is the second line of \eqref{gb:Lj} and the exponent $m+1$.
The $O(N^{-2})$ error in the lemma is small enough to give a
relative $O(n^{-1})$ error even in this second case.  Finally
\eqref{gb:t-asymptotic} implies
$\lim_{n\to\infty}|t_{n,j}|^{1/(2n)}=\lambda_j^{-1}$.
The Cauchy--Hadamard formula proves the exact radius.  This also
provides another proof of independence: the smallest surviving
radius in any nontrivial linear combination cannot cancel.
\end{proof}

\subsection{Eventual optimal reconstruction for every fixed gap}

Set
\begin{equation}\label{gb:asymptotic-normalization}
 D_m=m^2+\lfloor m/2\rfloor,\qquad
 \mathcal A_m=
 \frac{\prod_{j=1}^m|L_j|}{K_m}
       \left(\prod_{j=1}^m\xi_j\right)
       \prod_{j<k}(\xi_j-\xi_k)>0.
\end{equation}

\begin{theorem}[Fixed-gap determinant asymptotic]\label{gb:eventual}
For each fixed integer $m\ge1$, as $b\to\infty$,
\begin{equation}\label{gb:det-asymptotic}
 \boxed{\displaystyle
 \det F_{b,m}\sim
 (-1)^{\lfloor m^2/4\rfloor}\mathcal A_m
 (2b)^{-D_m}\bigl[\pi^{2m}(m!)^2\bigr]^{-b}.}
\end{equation}
In particular, there exists an integer $B(m)>m$ such that every
pair $(p,q)=(2b-2m,2b)$ with $b\ge B(m)$ has an optimal first-$q$
formal moment prefix.  Its late-block determinant has eventual sign
\begin{equation}\label{gb:eventual-sign}
 \operatorname{sgn}\det M_{b-m,b}
        =(-1)^{b(b-1)/2+\lfloor m^2/4\rfloor}.
\end{equation}
\end{theorem}

\begin{proof}
From column $j$ of $(t_{b+i,j})_{i,j=1}^m$, factor
$L_j\lambda_j^{-2b}(2b)^{-d_j}$.  By
\eqref{gb:t-asymptotic}, the remaining matrix tends entrywise to
$V=(\xi_j^i)_{i,j=1}^m$.  Its determinant is
\[
 \det V=\left(\prod_j\xi_j\right)
                   \prod_{j<k}(\xi_k-\xi_j)\ne0.
\]
The distinct numbers $\xi_1>\cdots>\xi_m>0$ give sign
$(-1)^{m(m-1)/2}$, which cancels the same sign in
\eqref{gb:exact-Gamma-determinant}.  Furthermore
\[
 \prod_j\operatorname{sgn}L_j
    =(-1)^{\sum_{j=1}^m\lfloor j/2\rfloor}
    =(-1)^{\lfloor m^2/4\rfloor},\qquad
 \sum_jd_j=m^2+\lfloor m/2\rfloor,
\]
and
$\prod_j\lambda_j^{-2b}=[\pi^{2m}(m!)^2]^{-b}$.
These identities prove \eqref{gb:det-asymptotic} with the positive
constant \eqref{gb:asymptotic-normalization}.  A sequence with a
nonzero asymptotic equivalent is nonzero and has the same sign
for all sufficiently large indices.  Apply the exact positive-factor
bridge \eqref{gb:bridge-formula} and then \eqref{gb:full-prefix}.
\end{proof}

\begin{remark}[Scope of the theorem]\label{gb:scope}
The assertion is eventual for each \emph{fixed gap} $2m$ and
does not provide a numerical value of $B(m)$.  It leaves only
finitely many possible exceptional pairs on each such ray, but
does not prove the entire gap-six family or the general
all-unequal-even conjecture.  No uniform assertion as $m$ grows
with $b$ follows from the fixed-$m$ errors used here.
\end{remark}

For $m=1$, \eqref{gb:evaluation} gives $K_1=8\pi$, and
\eqref{gb:asymptotic-normalization} gives $\mathcal A_1=4/\pi$.
The fundamental adjacent-gap sequence therefore satisfies
\[
        C^{[1]}_{b+1,1}\sim\frac{2}{\pi b}\,\pi^{-2b}.
\]
For $m=2$ the eventual sign of $\det F_{b,2}$ is negative,
consistent with the inherited theorem for every gap-four pair.
For $m=3$ it is positive.  Finite exact checks give
$\det F_{b,3}<0$ at $b=4,5,6$ and $\det F_{b,3}>0$ at
$7\le b\le40$.  These observations illustrate a sign transition;
the proof of eventual positivity is \cref{gb:eventual}, and the
finite checks are not a proof of nonvanishing for every gap-six pair.

\subsection{All-orders coefficients and a harmonic-number correction}

The Gamma construction supplies more than a leading equivalent.
It gives an explicit Poincare expansion to every finite order,
and the first determinant correction has a concise harmonic-number
form.  Let Euler numbers be defined by
\[
                \operatorname{sech}z=
                \sum_{k\ge0}E_{2k}\frac{z^{2k}}{(2k)!}.
\]
Define rational numbers $w_{m,r}$ by the formal identity
\begin{equation}\label{gb:w-series}
 \sum_{r\ge0}w_{m,r}X^r
   =\exp\!\left(\sum_{k\ge1}
            \frac{m|E_{2k}|}{k\,2^{2k+1}}X^k\right).
\end{equation}
Thus $w_{m,0}=1$ and $w_{m,1}=m/8$.
For $\ell\ge0$ set
\begin{equation}\label{gb:b-coefficients}
 b_{m,\ell}=\sum_{r=0}^{\lfloor\ell/2\rfloor}
  w_{m,r}\frac{(m+2r)_{\ell-2r}}{(\ell-2r)!}
                     (i/2)^{\ell-2r}.
\end{equation}

\begin{proposition}[All-orders mode construction]\label{gb:all-orders}
For each fixed mode $j$, with $\lambda_j,a_j$ as in
\eqref{gb:mode-constants},
\begin{equation}\label{gb:all-orders-moments}
 \int_0^\infty(x+i/2)^NP_j(x)W_m(x)\dd x
 \sim a_je^{i\lambda_j/2}
       \sum_{\ell\ge0}b_{m,\ell}
       \frac{\Gamma(N-m-\ell+1)}{
                            \lambda_j^{N-m-\ell+1}}.
\end{equation}
This is a Poincare expansion: truncation after any fixed $\ell=L$
has error of size
$O(\Gamma(N-m-L)\lambda_j^{-N+m+L})$.
Taking imaginary parts retains even $\ell$ when $j$ is odd and
odd $\ell$ when $j$ is even.  In particular,
\begin{equation}\label{gb:t-first-correction}
 t_{n,j}=L_j\lambda_j^{-2n}(2n)^{-d_j}
       \left(1+\frac{d_j(d_j-1)}{4n}+O(n^{-2})\right).
\end{equation}
\end{proposition}

\begin{proof}
The shifted logarithmic Stirling expansion is
\[
 \log\Gamma(z+\alpha)\sim
 (z+\alpha-\tfrac12)\log z-z+\tfrac12\log(2\pi)
 +\sum_{n\ge1}\frac{(-1)^{n+1}B_{n+1}(\alpha)}{n(n+1)z^n},
\]
with a remainder of the next prescribed algebraic order in each
closed subsector of $|\arg z|<\pi$ \cite{dlmf-stirling}.
Apply it with $\alpha=1/4$ and $z=\pm ix/2$.
The conjugate pair cancels odd inverse powers.  Using
\[
 B_{2k+1}(1/4)=-\frac{(2k+1)E_{2k}}{4^{2k+1}}
\]
gives
\[
 \log W_m(x)\sim
 \log A_m^*-m\log x-m\pi x
       +\sum_{k\ge1}\frac{m|E_{2k}|}{k\,2^{2k+1}x^{2k}}.
\]
Exponentiation to any fixed order proves
\begin{equation}\label{gb:W-all-orders}
 W_m(x)\sim A_m^*x^{-m}e^{-m\pi x}
                       \sum_{r\ge0}w_{m,r}x^{-2r}.
\end{equation}
The corrections to the leading exponential in $P_j$ are themselves
exponentially smaller, so they contribute less than any fixed
inverse power after the leading density has been removed.

For each term $x^{-m-2r}e^{-\lambda_jx}$ in the resulting density,
repeat the shift $y=x+i/2$ from \cref{gb:Laplace}.  Its algebraic
factor becomes
\[
 y^{-m-2r}\left(1-\frac{i}{2y}\right)^{-m-2r}
   =y^{-m-2r}\sum_{h\ge0}
       \frac{(m+2r)_h}{h!}\left(\frac{i}{2y}\right)^h.
\]
Collecting terms with $2r+h=\ell$ gives exactly
\eqref{gb:b-coefficients}.  To justify any fixed truncation, split
at fixed $X>1$ as in the lemma, retain the finitely many necessary
density terms, and bound its real remainder by the first omitted
power of $x^{-1}$ times the leading exponential.  The split at
$N^{2/3}$ bounds its moment by the corresponding Gamma ratio.
The finitely many model contours have the same pole-free deformation,
and Taylor's remainder for $(1-i/(2y))^{-m-2r}$ is uniform for
$y\ge X$.  These bounds prove the stated remainder in
\eqref{gb:all-orders-moments}; no contour continuation of the real
Stirling remainder is used.

For even $\ell$ the coefficient $b_{m,\ell}$ is real, and for
odd $\ell$ it is purely imaginary.  Thus the phase
$e^{ij\pi/2}$ gives the asserted parity selection.  At odd $j$
the leading imaginary term has $\ell=0$ and the next possible
one has $\ell=2$.  At even $j$ the leading one has $\ell=1$,
with $b_{m,1}=im/2$, and the next possible one has $\ell=3$.
In either case the first relative inverse-power correction comes
solely from the factorial ratio
\[
 \frac{\Gamma(N-d+1)}{\Gamma(N+1)}
 =N^{-d}\left(1+\frac{d(d-1)}{2N}+O(N^{-2})\right)
 \qquad(d\in\mathbb N).
\]
For integer $d$ this follows directly by expanding
$\prod_{r=0}^{d-1}(N-r)^{-1}$.  Set $N=2n$ and $d=d_j$
to obtain \eqref{gb:t-first-correction}.
\end{proof}

\begin{theorem}[The harmonic first correction]\label{gb:harmonic-correction}
Let $e=\lfloor m/2\rfloor$ and $H_0=0$.  Define
\begin{equation}\label{gb:kappa}
 \boxed{\displaystyle
 \kappa_m=-\frac{m^2(m+3)+e(2m+3)}4
       +\sum_{r=1}^{e}r\bigl(H_{m+2r}-H_{m-2r}\bigr).}
\end{equation}
Then
\begin{equation}\label{gb:det-first-correction}
\begin{split}
 \det F_{b,m}={}&(-1)^{\lfloor m^2/4\rfloor}\mathcal A_m
       (2b)^{-D_m}\bigl[\pi^{2m}(m!)^2\bigr]^{-b}\\
 &\times\left(1+\frac{\kappa_m}{b}+O(b^{-2})\right).
\end{split}
\end{equation}
More generally, this determinant has an expansion to every finite
inverse-power order.  For fixed $m$, all its normalized coefficients
belong to $\Q[\pi^2]$.
\end{theorem}

\begin{proof}
Factor the same column scales as in the proof of
\cref{gb:eventual}.  By \eqref{gb:t-first-correction}, the remaining
entry in row $i$, column $j$ is
\[
 \xi_j^i\left\{1+\frac1b
           \left(\frac{d_j(d_j-1)}4-id_j\right)+O(b^{-2})\right\}.
\]
For $V=(\xi_j^i)_{i,j=1}^m$, the first relative correction in its
determinant is therefore
\begin{equation}\label{gb:kappa-trace}
 \frac14\sum_{j=1}^md_j(d_j-1)-\sum_{j=1}^md_j\eta_j,
 \qquad
 \eta_j=\left(V^{-1}\operatorname{diag}(1,\ldots,m)V\right)_{jj}.
\end{equation}
The diagonal quantities can be computed without matrix inversion.
Differentiating the Vandermonde product gives
\[
 \eta_j=\xi_j\frac{\partial}{\partial\xi_j}\log\det V
       =1+\sum_{k\ne j}\frac{\xi_j}{\xi_j-\xi_k}.
\]
For $\xi_j=(j\pi)^{-2}$, partial fractions give
\begin{equation}\label{gb:eta-harmonic}
 \eta_j=m+\frac34+\frac j2(H_{m-j}-H_{m+j}).
\end{equation}
Indeed $\xi_j/(\xi_j-\xi_k)=k^2/(k^2-j^2)$, and
\[
 \sum_{k\ne j}\left(\frac1{k-j}-\frac1{k+j}\right)
        =H_{m-j}-H_{m+j}+\frac{3}{2j}.
\]
Also $\sum_j\eta_j=m(m+1)/2$, either by the trace or by
pairing the summands in the Vandermonde expression.  Since
$d_j=m+\ind_{2\mid j}$, substitution into
\eqref{gb:kappa-trace} gives
\[
 -\frac{m^2(m+3)}4+\frac{me}{2}
                         -\sum_{r=1}^e\eta_{2r}.
\]
Equation \eqref{gb:eta-harmonic} now simplifies it to
\eqref{gb:kappa}, proving \eqref{gb:det-first-correction}.

For higher orders, insert any desired finite truncation of
\eqref{gb:all-orders-moments} into the fixed-size determinant.
Its leading Vandermonde is nonzero, so division by that leading
term gives a valid Poincare expansion.  Relative to $L_j$, parity
selection leaves only even powers of $\lambda_j=j\pi$ in the
mode corrections.  The factorial-ratio expansions have rational
coefficients, while ratios of Vandermonde expressions in the
$\xi_j$ cancel their common powers of $\pi^{-2}$.  Thus each
normalized determinant coefficient is a polynomial in $\pi^2$
with rational coefficients, as claimed.
\end{proof}

For orientation, the first correction values are
\[
 \kappa_1=-1,\qquad\kappa_2=-\frac{14}{3},\qquad
 \kappa_3=-\frac{217}{15},\qquad\kappa_4=-\frac{949}{35}.
\]
The series in \eqref{gb:w-series} and
\eqref{gb:all-orders-moments} are formal asymptotic constructions;
convergence is not asserted.  The argument also does not establish
a complete exponentially improved expansion or an explicit error
constant sufficient to certify a numerical $B(m)$.

\subsection{A hypergeometric form of the null equation}

The differential equation itself has a useful classical
hypergeometric form.  Let $Q$ be any even analytic null germ of
$\mathcal L_m$ and write $c=Q^{(2m)}(0)$.  On the local branches
near $u=1$ with $\log1=0$ and $\sqrt1=1$, define
\[
 Y(u)=\frac{Q(\tfrac12\log u)}{\sqrt u-1},\qquad
                    \Theta=u\frac{d}{du}.
\]
The apparent singularity at $u=1$ is removable because $Q(t)=O(t^2)$.
Then
\begin{equation}\label{gb:hypergeometric}
 \left[\Theta^{2m}-u(\Theta+\tfrac12)^{2m}\right]Y(u)
            =-\frac{c}{2^{2m+1}}(1+\sqrt u).
\end{equation}
To verify it, set $Q(t)=(e^t-1)F(t)$, so that
$Q(t)\coth(t/2)=(e^t+1)F(t)$.  Multiplication of
$\mathcal L_mQ=0$ by $e^t-1$ gives
\[
 D^{2m}F-e^{2t}(D+1)^{2m}F=-\tfrac c2(1+e^t).
\]
Now $u=e^{2t}$ and $D=2\Theta$, proving
\eqref{gb:hypergeometric}.  Its homogeneous operator is the standard
generalized hypergeometric operator associated with
\[
 {}_{2m}F_{2m-1}
 \left(\begin{matrix}\tfrac12,\ldots,\tfrac12\\1,\ldots,1\end{matrix};u\right)
\]
\cite{dlmf-hypergeometric}.  This operator identification is classical
algebra.  The Gamma basis selects and normalizes the relevant
analytic null germs at $u=1$ and connects them to the harmonic-tail
coefficient matrix; that connection is the substantive result here.

\subsection{Verification and the next finite obstruction}

The accompanying scripts separate exact algebra from numerical
diagnostics.  Direct Gamma quadrature at $45$ decimal digits for
$m=1,2,3$ and $1\le n\le m+3$ gives all 18 coefficient recurrence
residuals below $10^{-46}$ and all six top-derivative identity
residuals below $10^{-42}$, with the normalizations recorded
separately. A separate calculation evaluates $T_m$
by Gamma quadrature at $60$ digits, then propagates the exact
recurrence to compare the determinant with
\eqref{gb:det-asymptotic}.  The recorded ratios are
\begin{center}
\begin{tabular}{@{}crrrr@{}}
\toprule
 & $b=20$ & $b=40$ & $b=80$ & $b=120$\\
\midrule
$m=1$ & $0.951663$ & $0.975424$ & $0.987607$ & $0.991714$\\
$m=2$ & $0.796107$ & $0.891056$ & $0.943647$ & $0.961999$\\
$m=3$ & $0.456866$ & $0.687818$ & $0.832149$ & $0.885309$\\
\bottomrule
\end{tabular}
\end{center}
Their approach to one is a consistency check on the explicit
normalization; the infinite theorem is proved above independently
of these finite computations.

The most direct continuation is to make the remainder in
\eqref{gb:det-asymptotic} effective, obtain a certified cutoff for
each fixed $m$, and use exact rational arithmetic below it.  The
gap-six case is a natural first target because its initial sign
transition rules out a blanket positivity shortcut.  Further questions
include exact evaluations of $K_m$, a classification of finite
sign changes, and uniform asymptotics when $m$ grows with $b$.
The Gamma moments also suggest parameter deformations of
\eqref{gb:W} and \eqref{gb:cot-operator}; any such extension must
recheck the contour residues and the original coordinate bridge.

% End included source: sections/08_gamma_basis
% Begin included source: sections/09_audit_research
\section{Research status, verification, and next questions}
\label{sec:audit}

\subsection{What has been established}

The proof-bearing deliverable is the article itself. The programs
check finite algebra and independent numerical instances; they are
not substituted for its convergence, continuation, contour, or
asymptotic arguments. The principal advances can be mapped precisely
to the source questions.

\begin{longtable}{@{}>{\raggedright\arraybackslash}p{0.24\textwidth}>{\raggedright\arraybackslash}p{0.45\textwidth}>{\raggedright\arraybackslash}p{0.25\textwidth}@{}}
\toprule
Source question & Result supplied here & Remaining scope\\
\midrule
\endfirsthead
\toprule
Source question & Result supplied here & Remaining scope\\
\midrule
\endhead
Individual ordered mixed labels &
Theorem~\ref{thm:ordered} gives the exact two-particle identity on an
open complex domain. Theorems~\ref{thm:poles} and
\ref{thm:digamma-master} give its continuation and all positive-order
values. &
Higher rank is handled by finite localization, not a universal
single-zeta assertion.\\[5pt]
Noninteger cutoff generators &
Theorem~\ref{thm:cutoff} and \eqref{eq:b-beta-lerch} give a finite
polylogarithm and beta--Lerch transform, entire in both orders on
the disk. &
The endpoint requires its own convergence domain and collective
cancellations.\\[5pt]
Other annihilating stencils &
Section~\ref{sec:localization} proves a finite block formula for
critical integer stencils and a stronger reduction under one further
annihilation condition. &
This is a sufficient family, not a classification of arbitrary
noninteger grids.\\[5pt]
All unequal even harmonic-tail prefixes &
Section~\ref{gb:section} proves eventual optimality on every fixed-gap
ray, with exact Gamma normalization and a nonzero signed asymptotic. &
No effective threshold is provided. The full all-pairs conjecture,
including complete gap-six coverage, remains open here.\\[5pt]
Stieltjes derivatives and primitives &
Theorems~\ref{thm:subtracted-jets} and \ref{thm:beta-stieltjes}
provide ordinary all-order identities. Equations
\eqref{eq:gamma-primitive} and \eqref{eq:arcsine-primitive}
give compatible primitives. &
Elementary symmetric arcsine cases are recurrence consequences;
they are not independent priority claims.\\[5pt]
Short Gaussian $S_6$ and revised $S_8$ &
Their present source status has been checked and preserved. &
Neither identity is proved or refuted by this article.\\
\bottomrule
\end{longtable}

The incoming bilinear and tail-kernel reports already supply complete
smaller-order-ten and smaller-order-twelve reconstruction
theorems \cite{repo:bilinear,repo:tail-kernels}. They should
retain their existing attribution and audit status. Reproving or
relabeling an already supplied finite case would not constitute a
new research advance; the fixed-gap theorem instead changes which
infinite direction is controlled.

\subsection{Audit findings and a proposed correction}

No newly false inherited mathematical theorem was identified in the
material used for this article. There is a concrete editorial defect:
the manuscript README repeats consecutively the complete paragraph
describing the 1,032-input first-letter-$X$ family, its rank 1,008,
and its kernel dimension 24 \cite{repo:manuscript}. The package includes
a narrowly scoped patch removing the second copy. This is a proposed
README correction, not a modification of a mathematical statement.

Several distinctions matter when the new results are integrated.
The source already records that the old $S_8$ candidate was rejected
and the revised candidate is still conjectural. The small residual
enclosures for $S_6$ and revised $S_8$ establish proximity, not zero.
Likewise, the canonical modular relation searches do not assert
rational nonmembership or a counterexample to a period equality
\cite{repo:ledger,repo:gaussian}. These are inherited status safeguards,
not newly discovered errors.

For the present contribution, the analogous precise statements are:
the zero-weight limits in Section~\ref{sec:continuation} concern a
meromorphic continuation; the constant in \eqref{eq:hard-cutoff} is
not the fixed-coordinate Laurent constant; block formulas outside
their separated convergence domain are interpreted collectively;
and full rank of a formal moment matrix does not prove arithmetic
independence of the evaluated constants. The Gamma argument also
shows why positivity of the scalar recurrence does not imply
positivity of every minor: even modes lose their leading Laplace
phase, and the determinant has the parity-dependent sign derived
in Section~\ref{gb:section}.

\subsection{Reproducibility and the independent checks}

The package contains modular TeX, a self-contained assembled TeX
file, the compiled PDF, a build command, a claim-status file, exact
source locations, and portable verification programs. Source
repository material was inspected at the pinned revision and kept
unchanged. The generated correction patch is separate.

The ordered-boundary suite checks the finite summation-by-parts
identity on 240 arbitrary rational sequence/cutoff cases, the rational
partial-fraction identity on 36 order pairs, 570 parity coefficients,
420 coefficients of the cutoff generator, and all four entries of
the zero-weight directional table. Four explicit infinite-sum values
are also compared with directly accumulated original ordered sums,
using both positive tail envelopes and separately labelled accelerated
diagnostics.

The corollary suite checks 13 exact master-function derivatives and
special values and nine exact formal Stieltjes coefficients. It
compares the geometric master with the original resolvent-weighted
ordered sum, tests two subtracted harmonic sums, follows six
unaccelerated hard-cutoff limits, and evaluates nine arcsine integrals.
The beta--Stieltjes suite checks five Bell--Leibniz coefficients,
independent Stieltjes-value samples, and twelve numerical beta
integrals through jet order two, including independently varied
shifts. Its expansion and endpoint truncations are recorded
separately from floating-point roundoff.

The localization suite supplies 147 exact checks, including
finite-support transport, the full contact table, both rank-three
two-index sectors, singleton cancellation through rank five, and a
deliberately corrupted control. The Gamma suite compares directly
integrated basis coefficients with the rational recurrence for
$m=1,2,3$, checks the exact inhomogeneous moment, and records
finite-index ratios to the determinant asymptotic. A separate exact
suite checks the Bernoulli matrix against its complementary-minor
normalization and the harmonic correction. The independent proof
audit rederives the contour shifts, both residues, the initial
determinant, the parity-sensitive asymptotics, and the source bridge.
This is internal independent checking, not external peer review or
proof-assistant verification.

All numerical comparisons are diagnostics at their stated working
precisions. The proof of eventual nonvanishing is the nonzero
limiting Vandermonde determinant, not an extrapolation of the
displayed finite data. In particular, a finite gap-six sign table
does not provide the missing effective cutoff.

\subsection{Further research questions}

The following questions are arranged by how directly the present
proofs supply a route forward. They concern exact identities and
structural classification; estimates enter only where needed to
turn an eventual identity theorem into a finite decision.

\begin{enumerate}[label=\textbf{\arabic*.},leftmargin=*]
\item \textbf{Complete the fixed-gap prefix conjectures.}
For each fixed $m$, turn the nonzero asymptotic into an explicit
computable threshold $B(m)$ and combine it with exact rational
determinants below that threshold. The specific next case is gap
six. The exact values at $b=4,5,6$ are negative, and the checked
values $7\le b\le40$ are positive. The strengthened conjecture
\[
 \det F_{b,3}>0\qquad(b\ge7)
\]
is consistent with both the finite data and the proved eventual
sign, but is not proved here. A successful completion must justify
the transition from a finite table to all indices.

\item \textbf{Evaluate the Gamma normalization constants.}
The positive integral $K_m$ is exact, and $K_1=8\pi$.
Find closed Gamma, hypergeometric, or modular descriptions of
$K_2,K_3,\ldots$. The multiple-orthogonality interpretation suggests
studying the associated moment determinant rather than integrating
each entry separately. A recurrence in $m$, if it exists, would be
particularly useful.

\item \textbf{Connect the two eventual regimes.}
The new theorem fixes the gap while the earlier theorem fixes the
smaller order. Can one obtain a two-parameter determinant expansion
that covers the region between those rays? The limiting nodes
$(j\pi)^{-2}$ become less separated as $m$ grows, so a proof must
control the full determinant, not merely reuse fixed-$m$ error
constants. This is a possible route toward the all-pairs conjecture.

\item \textbf{Classify the finite sign changes.}
The exact Gamma representation is valid at every $b>m$, while its
sign is presently asymptotic. Determine the number and locations of
the finite sign transitions for each $m$. A variation-diminishing
statement adapted to the polynomials
$\operatorname{Im}(x+i/2)^{2n}$ could replace a false appeal to
unqualified total positivity.

\item \textbf{Deform the Gamma basis.}
Replace the quarter shift by a parameter and seek the corresponding
shift equation, residue forcing, and reconstruction matrix. Identify
which deformations preserve a finite set of sinh/tanh modes and
which create additional residue terms. Differentiating a valid
deformation would give new exact polygamma-weighted orthogonality
identities; the deformation itself must be proved before its
derivatives are used.

\item \textbf{Determine the actual higher-rank period coordinates.}
The block theorem gives at most
$\lfloor(r+1)/2\rfloor$ free ordered indices. At rank three the
two-index contact sector is nonzero, but this does not prove
arithmetical depth two. Determine which integer orders, shift
collisions, or label symmetries reduce these sectors to single
Hurwitz values, and which require genuine nested periods.
The explicit contact polynomial provides a finite starting point.

\item \textbf{Compute all-rank polar geometry from the blocks.}
The two-particle oblique residues are explicit polynomials.
At higher rank, the finite sector formula should organize polar
hyperplanes by free blocks and their mergers. Prove a complete
local Laurent description, including its behavior when the original
finite-difference weights vanish. This would extend the directional
table without confusing a pointwise specialization with a limit.

\item \textbf{Classify stencils by ordered boundary support.}
The annihilation conditions proved here are sufficient. Seek
necessary conditions for a prescribed maximum number of free
blocks, first for finite integer stencils. A discrete-residue
criterion might recognize additional cancellations invisible to
degree counting. Nonuniform or noninteger grids require a separate
analysis because there is no common integer transport lattice.

\item \textbf{Find new exact beta--Stieltjes evaluations.}
For rational $\sigma$, special shifts may combine the beta weight
with multiplication or reflection formulas. Determine when
\eqref{eq:beta-stieltjes-all} reduces to a finite expression in
zeta derivatives, log-Gamma values, and lower Stieltjes constants,
rather than the generally infinite but geometric expansion
\eqref{eq:beta-jet-zeta}. Confluent shift derivatives offer a
systematic way to introduce polygamma factors.

\item \textbf{Construct higher-rank cutoff formulas.}
Insert a common outer cutoff before transporting the stencils.
Track the factors of $1-z$ attached to free collision blocks and
determine whether the endpoint depth reduction can be read directly
from them. Section~\ref{sec:polylog} proves this mechanism at rank
two. A higher-rank formula should retain all finite endpoint terms
before any limiting operation.

\item \textbf{Return to Gaussian periods with a new functional route.}
The short $S_6$ and revised $S_8$ questions remain worthwhile, but
the inspected finite relation spaces have already been analyzed
carefully. New work should supply a functional identity, a genuinely
larger motivic or iterated-integral relation, or a rigorously changed
candidate. Further decimal agreement alone cannot settle either
equality. The ordered kernels here may suggest new generating
relations, but no such implication is presently proved.
\end{enumerate}

\subsection{Suggested integration}

The ordered identity and its cutoff generator naturally continue
the finite-difference report. The beta--Stieltjes identities belong
with the Hurwitz and Stieltjes calculus, with their ordinary
convergence hypotheses retained. The all-rank transport theorem
can be kept as a separate structural continuation. The Gamma basis
should follow the exact fixed-gap determinant reduction in the
harmonic-tail Gamma branch, with its inherited coordinate convention
printed explicitly.

The accompanying integration note lists these destinations, theorem
dependencies, a proposed status update, and the README patch.
Canonical acceptance should record the proved fixed-gap eventual
statement at its actual scope; it should not silently promote the
remaining full reconstruction or Gaussian conjectures.

% End included source: sections/09_audit_research
% Begin included source: references
\clearpage
\begin{thebibliography}{99}
\addcontentsline{toc}{section}{References}
\raggedright

\bibitem{repo:manuscript}
ProveIt Contributors,
\emph{Polylogarithms and their Arithmetic Bridges}, collective
manuscript, README and source chapters.
Repository snapshot \repo, 11 October 2026.
\href{https://github.com/VladimirReshetnikov/ProveIt/blob/e3d7ec5d3f9be259c1f91edaea4389f4e79b5b03/Analysis/Polylogarithms/docs/manuscript/README.md}{Pinned manuscript README}.

\bibitem{repo:ledger}
ProveIt Contributors,
\emph{Editorial Ledger}, canonical claim and integration statuses,
same repository snapshot.
\href{https://github.com/VladimirReshetnikov/ProveIt/blob/e3d7ec5d3f9be259c1f91edaea4389f4e79b5b03/Analysis/Polylogarithms/docs/manuscript/EDITORIAL-LEDGER.md}{Pinned editorial ledger}.

\bibitem{repo:gaussian}
ProveIt Contributors,
\emph{Revised \(S_8\) candidate}, manuscript Chapter 4 source,
with the surrounding Gaussian status and certification discussion.
\href{https://github.com/VladimirReshetnikov/ProveIt/blob/e3d7ec5d3f9be259c1f91edaea4389f4e79b5b03/Analysis/Polylogarithms/docs/manuscript/chapters/04-S8-candidate.tex}{Pinned candidate source}.

\bibitem{repo:finite}
ProveIt Contributors,
\emph{Finite-Difference Determinants}, incoming research report,
11 October 2026; especially Sections 2--4 and 9, and the subsequent
finite-algorithm continuation.
\href{https://github.com/VladimirReshetnikov/ProveIt/blob/e3d7ec5d3f9be259c1f91edaea4389f4e79b5b03/docs/incoming/ProveIt_Finite_Difference_Determinants_2026-10-11.zip}{Pinned incoming archive}.

\bibitem{repo-gamma}
ProveIt Contributors,
\emph{Harmonic-tail Gamma continuation}, Section 2,
\texttt{02\_gap\_four.tex}; inherited moment coordinates,
operator conjugacy, and Theorem \texttt{gap:complementary}.
\href{https://github.com/VladimirReshetnikov/ProveIt/blob/e3d7ec5d3f9be259c1f91edaea4389f4e79b5b03/Analysis/Polylogarithms/docs/reports/stieltjes-harmonic-resonance/continuations/harmonic-tail-gamma/sections/02_gap_four.tex}{Pinned source and exact determinant normalization}.
The earlier adjacent-tail conventions are in
\href{https://github.com/VladimirReshetnikov/ProveIt/blob/e3d7ec5d3f9be259c1f91edaea4389f4e79b5b03/Analysis/Polylogarithms/docs/reports/stieltjes-harmonic-resonance/continuations/differential-polygamma-lerch-jets/sections/08_adjacent_tails.tex}{the differential-polygamma continuation}.

\bibitem{repo:bilinear}
ProveIt Contributors,
\emph{Bilinear Tails and Harmonic Determinants}, incoming report,
11 October 2026; rational boundary certificates, cutoff jets,
and order-ten reconstruction.
\href{https://github.com/VladimirReshetnikov/ProveIt/blob/e3d7ec5d3f9be259c1f91edaea4389f4e79b5b03/docs/incoming/ProveIt_Bilinear_Tails_and_Determinants_2026-10-11.zip}{Pinned incoming archive}.

\bibitem{repo:tail-kernels}
ProveIt Contributors,
\emph{Tail Kernels and Boundary Products}, incoming report,
11 October 2026; Sections 6 and 7 supply complete smaller-order-ten
and smaller-order-twelve reconstruction theorems.
\href{https://github.com/VladimirReshetnikov/ProveIt/blob/e3d7ec5d3f9be259c1f91edaea4389f4e79b5b03/docs/incoming/ProveIt_Tail_Kernels_and_Boundary_Products_2026-10-11.zip}{Pinned incoming archive}.
All seven incoming archives and their specific roles are listed in
the accompanying source map.

\bibitem{repo:quartic}
ProveIt Contributors,
\emph{Quartic Spectral Reconstruction}, incoming report,
11 October 2026; Section 3 proves eventual reconstruction with the
smaller order fixed, and Section 8 poses the order-ten/twelve tasks
subsequently addressed in the other incoming reports.
\href{https://github.com/VladimirReshetnikov/ProveIt/blob/e3d7ec5d3f9be259c1f91edaea4389f4e79b5b03/docs/incoming/ProveIt_Quartic_Spectral_Reconstruction_2026-10-11.zip}{Pinned incoming archive}.

\bibitem{dlmf:hurwitz}
NIST Digital Library of Mathematical Functions,
\emph{Hurwitz Zeta Function}, Section 25.11.
\url{https://dlmf.nist.gov/25.11}.
Definitions, continuation, argument derivatives, Gamma relation,
and asymptotic expansions.

\bibitem{dlmf:lerch}
NIST Digital Library of Mathematical Functions,
\emph{Lerch's Transcendent}, Section 25.14.
\url{https://dlmf.nist.gov/25.14}.

\bibitem{dlmf:gamma}
NIST Digital Library of Mathematical Functions,
\emph{Gamma Function: Functional Relations and Integral Representations},
Sections 5.5 and 5.9.
\url{https://dlmf.nist.gov/5.5};
\url{https://dlmf.nist.gov/5.9}.

\bibitem{dlmf-stirling}
NIST Digital Library of Mathematical Functions,
\emph{Gamma Function: Asymptotic Expansions}, Section 5.11.
\url{https://dlmf.nist.gov/5.11}.

\bibitem{dlmf-hypergeometric}
NIST Digital Library of Mathematical Functions,
\emph{Generalized Hypergeometric Functions: Differential Equations},
Section 16.8.
\url{https://dlmf.nist.gov/16.8}.

\bibitem{chen-singer}
S. Chen and M. F. Singer,
\emph{Residues and Telescopers for Rational Functions},
arXiv:1203.4200 (2012).
\url{https://arxiv.org/abs/1203.4200}.

\bibitem{arreche-sitaula}
C. E. Arreche and H. P. Sitaula,
\emph{A computational approach to rational summability and its
applications via discrete residues},
arXiv:2503.15636 (2025).
\url{https://arxiv.org/abs/2503.15636}.

\bibitem{matsumoto-ono-wakabayashi}
K. Matsumoto, T. Onozuka, and I. Wakabayashi,
\emph{Laurent series expansions of multiple zeta-functions of
Euler--Zagier type at integer points},
arXiv:1601.05918 (2016).
\url{https://arxiv.org/abs/1601.05918}.

\bibitem{shpot-paris}
M. A. Shpot and R. B. Paris,
\emph{Integrals of products of Hurwitz zeta functions via Feynman
parametrization and two double sums of Riemann zeta functions},
arXiv:1609.05658 (2016).
\url{https://arxiv.org/abs/1609.05658}.

\bibitem{ramanujan-transform}
D. Chakrabarti and G. K. Srinivasan,
\emph{On a remarkable formula of Ramanujan},
arXiv:1206.5122 (2012).
\url{https://arxiv.org/abs/1206.5122}.
Classical Gamma Fourier-transform context; no identification with the
specific fixed-gap basis proved here is asserted.

\end{thebibliography}

% End included source: references
\end{document}
